\documentclass[ecta,nameyear,final]{econsocart}
\RequirePackage[colorlinks,citecolor=blue,linkcolor=blue,urlcolor=blue]{hyperref}
\usepackage{amssymb,bm,booktabs,array,multirow,graphicx}
\usepackage{pgfplots}
\pgfplotsset{compat=1.18}
\usepgfplotslibrary{groupplots}

\startlocaldefs
\allowdisplaybreaks
\renewcommand{\topfraction}{0.9}\renewcommand{\bottomfraction}{0.7}\renewcommand{\textfraction}{0.08}\renewcommand{\floatpagefraction}{0.75}

\newcommand{\E}{{\mathbb E}}
\newcommand{\PP}{{\mathbb P}}
\newcommand{\dr}[1]{{\mathrm d}#1}
\newcommand{\iu}{{\mathrm i}}
\newcommand{\Rea}{\operatorname{Re}}
\newcommand{\Ima}{\operatorname{Im}}
\newcommand{\Var}{\operatorname{Var}}
\newcommand{\Cov}{\operatorname{Cov}}
\newcommand{\erf}{\operatorname{erf}}
\newcommand{\erfc}{\operatorname{erfc}}
\newcommand{\tr}{\operatorname{tr}}
\newcommand{\adj}{\operatorname{adj}}
\newcommand{\sech}{\operatorname{sech}}
\newcommand{\diag}{\operatorname{diag}}
\newcommand{\phih}{\hat\phi}
\newcommand{\tauh}{\hat\tau}
\newcommand{\Fk}{\textsf{\upshape DK}}
\newcommand{\Fu}{\textsf{\upshape DU}}
\newcommand{\Ft}{\textsf{\upshape DT}}
\newcommand{\Sk}{\textsf{\upshape SK}}
\newcommand{\Su}{\textsf{\upshape SU}}
\newcommand{\St}{\textsf{\upshape ST}}
\newcommand{\cF}{\mathcal{F}}
\newcommand{\cP}{\mathcal{P}}
\newcommand{\cQ}{\mathcal{Q}}
\newcommand{\cG}{\mathcal{G}}
\newcommand{\bX}{\boldsymbol{X}}
\newcommand{\bQ}{\boldsymbol{Q}}
\newcommand{\bK}{\boldsymbol{K}}
\newcommand{\bI}{\boldsymbol{I}}
\newcommand{\bGam}{\boldsymbol{\Gamma}}
\newcommand{\bSig}{\boldsymbol{\Sigma}}
\newcommand{\bP}{\boldsymbol{P}}
\newcommand{\bG}{\boldsymbol{G}}

\theoremstyle{plain}
\newtheorem{theorem}{Theorem}
\newtheorem{proposition}{Proposition}
\newtheorem{lemma}{Lemma}
\newtheorem{corollary}{Corollary}
\theoremstyle{definition}
\newtheorem{remark}{Remark}
\newcommand{\runin}[1]{\par\medskip\noindent\textit{#1.}\ \ignorespaces}

\endlocaldefs

\begin{document}

\begin{frontmatter}
\title{The Exact Limiting Distribution of the Dickey--Fuller $t$-Ratio}
\runtitle{The Limiting Distribution of the Dickey--Fuller $t$-Ratio}

\begin{aug}
\author[purdue]{\fnms{Yong}~\snm{Bao}\ead[label=e1]{ybao@purdue.edu}}
\author[hslu]{\fnms{Simon A.}~\snm{Broda}\ead[label=e2]{simon.broda@hslu.ch}}
\author[rotman,hkust]{\fnms{Raymond}~\snm{Kan}\ead[label=e3]{Raymond.Kan@rotman.utoronto.ca}}
\address[purdue]{\orgdiv{Department of Economics}, \orgname{Purdue University}}
\address[hslu]{\orgdiv{Institute of Financial Services IFZ}, \orgname{Lucerne University of Applied Sciences and Arts}}
\address[rotman]{\orgdiv{Rotman School of Management}, \orgname{University of Toronto}}
\address[hkust]{\orgdiv{School of Business and Management}, \orgname{Hong Kong University of Science and Technology}}
\end{aug}

\begin{abstract}
The limiting distribution of the Dickey--Fuller $t$-ratio, the reference distribution of the most widely used test in time series econometrics, has been tabulated by simulation since 1979.  We obtain its density and distribution function in closed form up to one integral, for the regression without deterministic terms, with an intercept, and with an intercept and a linear trend, under local alternatives $\phi=1+c/n$ on either side of the unit root, for the $t$-ratio centered at the unit root or at any local null value, and with a zero, a fixed or a stationary initial observation of the order of $\sqrt n$.  The joint moment generating function of the numerator and the denominator of the limit is a Fredholm determinant that is a polynomial of degree at most two in the argument of the numerator, because the numerator is a quadratic form of rank at most two in a few Gaussian functionals of the path, and the initial observation enters it through a closed-form quadratic term.  An inversion formula for ratios of the form $X/\sqrt Y$ then gives the density and the distribution function as integrals of special functions along a line in the complex plane, evaluated to eleven or twelve digits in a fraction of a second, and the moments of the limits are one-dimensional integrals of the moment generating function at real arguments.  The limiting power of the three tests, with a zero, a fixed or a stationary start, is then exact, and so are the power envelopes, the local power of the GLS-detrended tests, and confidence sets for the persistence parameter.
\end{abstract}

\begin{keyword}
\kwd{Unit root}
\kwd{Dickey--Fuller test}
\kwd{$t$-ratio}
\kwd{local-to-unity asymptotics}
\kwd{Fredholm determinant}
\kwd{inversion formula}
\kwd{parabolic cylinder function}
\kwd{power envelope}
\kwd{confidence sets}
\end{keyword}

\end{frontmatter}

\section{Introduction}
\label{sec:intro}

The Dickey--Fuller $t$-ratio is the least squares $t$-statistic for $\phi=1$ in the first-order autoregression $y_t=\alpha+\phi y_{t-1}+\sigma e_t$, computed from the regression of $y_t$ on $y_{t-1}$ (the statistic $\hat\tau$), on a constant and $y_{t-1}$ ($\hat\tau_\mu$), or on a constant, a linear trend and $y_{t-1}$ ($\hat\tau_\tau$) \citep{DickeyFuller1979,DickeyFuller1981,Fuller1996}.  Under the unit root its limiting distribution is a functional of Brownian motion \citep{DickeyFuller1979,ChanWei1987,Phillips1987}, and under the local alternatives $\phi=1+c/n$ a functional of an Ornstein--Uhlenbeck process.  Every critical value in use---the tables of \cite{Fuller1996}, the response surfaces of \cite{MacKinnon1996,MacKinnon2010} and \cite{CheungLai1995}, and the tables built into current software---comes from simulation, and so do the local power curves of the three tests \citep{Elliott1999,MullerElliott2003,HarveyLeybourneTaylor2009}.

The reason is that the $t$-ratio is not a ratio of quadratic forms.  For the coefficient statistic $n(\phih-1)$, whose limit is the ratio of two quadratic forms in the path, the joint characteristic function of numerator and denominator is a Fredholm determinant \citep{NabeyaTanaka1990}, and the distribution function follows by the inversion formula of \cite{Imhof1961} and \cite{GilPelaez1951} as a one-dimensional integral; \cite{NabeyaTanaka1990b} and \cite{NabeyaSorensen1994} tabulated the limiting distributions of the coefficient statistics in this way, under local alternatives and with a nonzero initial value, a drift and a trend, and \cite{Tanaka1990} and \cite{Perron1991b} treated the fixed and the stationary initial observation.  The limit of the $t$-ratio is the ratio of a quadratic form to the square root of another, which this approach does not reach \citep[Section~8.4]{Tanaka2017}.  What is known exactly is partial.  \cite{Abadir1992} computed the limiting density and distribution function of $\hat\tau$ numerically, from a formula with a double sum and an integral, and \cite{Abadir1995} derived a convergent series for its lower tail and one-dimensional integrals of parabolic cylinder functions for its upper tail, for $\hat\tau$ at the unit root with a zero start (\cite{Abadir1993a} did the same for the coefficient statistic).  \cite{Larsson1995} gave an exact formula for the upper tail of the limiting distribution of $\hat\tau$ under local alternatives and saddlepoint approximations to both tails, and \cite{Larsson1998} saddlepoint approximations to its lower tail; \cite{HitomiEtAl2022} wrote the limiting distributions of the coefficient and $t$ statistics without an intercept, with a zero or a fixed initial value, as double integrals of the joint density of a squared Bessel process and its time integral, whose integrands are triple sums of parabolic cylinder functions.  \cite{Nabeya1999} computed the first six moments of the three limiting null distributions, by the formula of \cite{Sawa1978}, and \cite{Larsson1997} the second moment of $\hat\tau_\mu$.  For $\hat\tau_\mu$ and $\hat\tau_\tau$, under the null and under local alternatives, no exact result for the distribution has been available, and the effect of the initial observation on the local power of the three $t$-tests has been studied by simulating the limiting functionals \citep{Elliott1999,MullerElliott2003,HarveyLeybourneTaylor2009}.

This paper gives the exact limiting density and distribution function of the three $t$-ratios, under the unit root and under local alternatives on either side of it, for the statistic centered at the unit root or at any local null value $\phi_0=1+c_0/n$, and with a zero, a fixed or a stationary initial observation of the order of $\sqrt n$, as one-dimensional integrals of closed-form integrands, and evaluates them to eleven or twelve digits in a fraction of a second.  The representation rests on two facts.  First, the limiting numerator and denominator of the $t$-ratio, $U=\int_0^1\tilde J\dr{J}$ and $V=\int_0^1\tilde J^2\dr{r}$ in the notation of Section~\ref{sec:model}, have a joint moment generating function that is a Fredholm determinant \citep{NabeyaTanaka1990,NabeyaSorensen1994}, and this determinant is a polynomial of degree one or two in the argument of $U$, because $U$ is a quadratic form of rank one or two in a few Gaussian functionals of the path, $J(0)$, $J(1)$, $\int_0^1J\dr{r}$ and $\int_0^1rJ\dr{r}$; a fixed initial observation adds to its logarithm a closed-form term that is quadratic in the start, and a stationary start, averaged over, leaves a determinant of the same degree.  Second, the distribution function of a ratio $X/\sqrt Y$ is, written out, a quadruple integral: two inversions of the joint moment generating function of $(X,Y)$ for the joint density, one over the parabola on which the ratio takes the value $q$, and one over $q$.  The last two act on the exponential alone and are done in closed form, which leaves a double inversion with an explicit kernel (Section~\ref{sec:inv}); the polynomial structure then makes the inversion in the argument of $U$ explicit, a scaled $\chi^2_1$ density for $\hat\tau$ and the density of a difference of two scaled $\chi^2_1$ variables, a Bessel function $K_0$, for $\hat\tau_\mu$, $\hat\tau_\tau$ and the three stationary cases, and what remains is one integral, over the argument of $V$, of parabolic cylinder functions and of Gaussian transforms of $K_0$ along a line in the complex plane.  The integrand decays only like $\exp(-\frac12|\Ima t|^{1/2})$ along the line and, in the lower tail, oscillates with a period proportional to the square of the argument; a partition of the integrand into an analytic part and an oscillating part, each integrated by a rule suited to it, gives twelve digits with a few hundred nodes.  A second representation of the same double integral, taken in the other order, serves as an independent check; it is in the Supplementary Material, with the quadrature.  The moments of the limits need no inversion at all: a gamma integral makes them one-dimensional integrals of the moment generating function at real arguments.

The principle is more general than the problem.  Whenever a statistic is a ratio $X/\sqrt Y$ in which $X+a$ is a quadratic form of low rank $r$ in a Gaussian vector and $Y$ a quadratic form in the same vector, the joint moment generating function is $e^{-as}$ times the inverse square root of a polynomial of degree $r$ in $s$ whose coefficients depend on $t$ alone, because under the tilt $e^{tY}$ the vector remains Gaussian and $X+a$ remains a quadratic form of rank $r$.  The inversion in $s$ is then the density of a weighted sum of $r$ squared normals, in closed form for $r\le2$, and the distribution of the ratio is one integral in $t$ of that density against an explicit Gaussian kernel.  The Dickey--Fuller limits are the case $r\le2$ with $a=\frac12$, and so are the limits of the unit root tests with a drift or a trend in the data-generating process.  The limit of the $t$-ratio in a predictive regression with a near-integrated regressor is the same ratio, scaled by the correlation of the innovations, convolved with an independent normal, since the part of its numerator that is orthogonal to the regressor's innovations is conditionally Gaussian given the path.

The limiting power of the three tests with a fixed or a stationary start, studied so far by simulation, is then exact: the start raises power, most for the test without deterministic terms, and lowers it slightly for the test with a trend close to the unit root.  The finite-sample distribution of the same statistics, which the limits approximate, is the subject of the companion paper \citep{BaoBrodaKan2026}, whose exact finite-sample powers we set beside the limits.  Two applications follow.  The coefficient statistic $n(\hat\phi-1)$ has the same limiting numerator and denominator, so that its local power, and that of the GLS-detrended tests of \cite{ElliottRothenbergStock1996}, come from the same formulas, exactly, with the power envelopes beside them; and the $t$-ratio centered at a local null value is the Wald statistic for that value, whose exact limit gives confidence sets for the persistence parameter by test inversion, with the initial condition.

Section~\ref{sec:model} sets out the models and their limits, Section~\ref{sec:mgf} the joint moment generating function, its structure and the initial observation, Section~\ref{sec:inv} the inversion formulas and the representation of the density and the distribution function, Section~\ref{sec:mom} the moments, Section~\ref{sec:results} the numerical results, and Sections~\ref{sec:app} and~\ref{sec:ci} the two applications.  Proofs are in Appendix~\ref{app:proofs}.  The Supplementary Material gives a second representation, the quadrature, the accuracy and the cost of the programs, and the checks against the closed forms, the moments, the critical values and the simulations in the literature.  The programs that produce every number in the paper accompany it, in MATLAB, Julia and C, the latter with a MATLAB interface.

\section{The $t$-ratios and their limits}
\label{sec:model}

Consider the first-order autoregression
\begin{equation}
y_t=\alpha+\delta t+\phi y_{t-1}+\sigma e_t,\qquad t=1,\ldots,n,
\label{eq:ar1}
\end{equation}
with $e_1,\ldots,e_n$ independent $N(0,1)$ and $\sigma>0$.  Following \cite{BaoKan2026} we label the cases by the initial observation, \textsf{D}eterministic, that is fixed, or \textsf{S}tationary, and by the deterministic terms, intercept \textsf{K}nown to be zero, intercept \textsf{U}nknown, or intercept and \textsf{T}rend unknown.  In case \Fk\ ($\alpha=\delta=0$ known, $y_0$ fixed) $\phi$ is estimated by the regression of $y_t$ on $y_{t-1}$, in case \Fu\ ($\delta=0$ known) by the regression on a constant and $y_{t-1}$, and in case \Ft\ by the regression on a constant, $t$ and $y_{t-1}$.  Let $\tauh$ be the least squares $t$-ratio for the hypothesis $\phi=1$, that is $\phih-1$ divided by its standard error: the Dickey--Fuller statistics $\hat\tau$, $\hat\tau_\mu$ and $\hat\tau_\tau$.  More generally let $\tauh_0$ be the $t$-ratio for the hypothesis $\phi=\phi_0$ with $\phi_0=1+c_0/n$, the Dickey--Fuller statistic at $c_0=0$ and the $t$-ratio centered at the true value at $c_0=c$, whose distribution is needed when a confidence interval for $\phi$ is obtained by inverting the test.  Cases \Sk, \Su\ and \St\ are cases \Fk, \Fu\ and \Ft\ with $|\phi|<1$ and the initial observation drawn from the stationary distribution of the process.

Let $\phi=1+c/n$ with $c$ fixed as $n\to\infty$, no drift, and $y_0=\sigma\sqrt n\,\gamma_0$.  Let $W$ be a standard Brownian motion on $[0,1]$ and
\begin{equation}
J(r)=\gamma_0e^{cr}+\int_0^re^{c(r-u)}\dr{W(u)},\qquad\dr{J}=cJ\dr{r}+\dr{W},
\label{eq:J}
\end{equation}
the Ornstein--Uhlenbeck process started at $\gamma_0$, where $\gamma_0$ is a constant in the fixed-start cases ($\gamma_0=0$ is the standard setting, in which the start is negligible against $\sqrt n$) and $\gamma_0\sim N\{0,-1/(2c)\}$ independent of $W$, the limit of the stationary distribution of $y_0/(\sigma\sqrt n)$, in the stationary cases, which require $c<0$.  Following \cite{MullerElliott2003} we also write $\gamma_0=\xi/\sqrt{-2c}$, so that $\xi$ is the initial deviation in units of the stationary standard deviation and the stationary start is $\xi\sim N(0,1)$.  Let $\tilde J$ be $J$ (\Fk, \Sk), $J$ less its mean $\int_0^1J\dr{r}$ (\Fu, \Su), or $J$ less its projection on $(1,r)$ in $L^2[0,1]$ (\Ft, \St).  Then, as $n\to\infty$ \citep{Phillips1987,ChanWei1987,NabeyaSorensen1994},
\begin{equation}
\tauh_0\Rightarrow\tau=\frac{U-c_0V}{\sqrt V},\qquad U=\int_0^1\tilde J(r)\dr{J(r)},\qquad
V=\int_0^1\tilde J(r)^2\dr{r},
\label{eq:limit}
\end{equation}
since $\tauh_0=\tauh-c_0/(n\,\mathrm{se})$ and $n\,\mathrm{se}\to V^{-1/2}$, where $\mathrm{se}$ is the standard error of $\phih$.  Here $U$, the limit of $n(\phih-1)$ times the limit of $n^2\mathrm{se}^2$, is defined by the formula $\int_0^1J\dr{W}+c\int_0^1J^2\dr{r}$ less the projection terms, since $\tilde J$ depends on the whole path and the stochastic integral is not an It\^o integral of $\tilde J$.  In cases \Fu, \Ft, \Su\ and \St\ the regression removes $\gamma_0$ itself and only $\gamma_0(e^{cr}-1)$ remains, which vanishes at $c=0$: the null limits do not depend on the start, and the start matters only under the alternative.  We write $\cF$ and $f$ for the distribution function and the density of the limit in question.

Two identities organize what follows.  By It\^o's formula, $U+\frac12$ is a quadratic form of rank one or two in a few Gaussian functionals of the path.  With $\bar J=\int_0^1J\dr{r}$ and $K=\int_0^1rJ\dr{r}$,
\begin{equation}
U+\frac12=\frac{J(1)^2-J(0)^2}{2}-L,\qquad
L=\begin{cases}
0 & (\textsf{K}),\\
\{J(1)-J(0)\}\bar J & (\textsf{U}),\\
\{J(1)-J(0)\}(4\bar J-6K)\\ \quad+6\{J(1)-\bar J\}(2K-\bar J) & (\textsf{T}),
\end{cases}
\label{eq:U}
\end{equation}
which for a zero start reduces to $\frac12J(1)^2$, $\frac12J(1)^2-J(1)\bar J$ and $\frac12J(1)^2-J(1)\bar J-12\{\frac12J(1)-\bar J\}(K-\frac12\bar J)$.  The rank is one in case \Fk\ and two otherwise; in the cases with deterministic terms $U$ is invariant to $J\mapsto J+a$ (\textsf{U}) or $J\mapsto J+a+br$ (\textsf{T}), as the regression is, so that $U+\frac12$ is a quadratic form in two invariant linear combinations of these functionals, $J(1)-J(0)$ and $J(1)+J(0)-2\bar J$ in case \textsf{U}.  And given $V$ the event $\{\tau\le x\}=\{U\le x\sqrt V+c_0V\}$ is linear in $U$.  The first identity makes the joint moment generating function of $(U,V)$ a polynomial in the argument of $U$ (Section~\ref{sec:mgf}); the second is what the inversion formulas of Section~\ref{sec:inv} exploit.  We write $R=U+\frac12$ where the quadratic form is in view, and $a(v)=\frac12+x\sqrt v+c_0v$, so that $\{\tau\le x\}=\{R\le a(V)\}$.

\section{The joint moment generating function}
\label{sec:mgf}

Let $M(s,t)=\E\exp(sU+tV)$, which is finite on a tube $\{(\Rea s,\Rea t)\in\mathcal{D}\}$ with $\mathcal{D}$ a convex neighborhood of the origin.  With a zero start it is known in closed form: the Fredholm determinants of \citet[Theorem~4]{NabeyaTanaka1990} for their Models~1, 2 and 4 give it in the three cases, after their $c$ is replaced by $-c$ and their statistic $T(\hat\rho-1)$ is identified with $U/V$; case \Fk\ is the moment generating function of \cite{White1958} and \cite{Phillips1987}.  We write the result in a form that displays its structure.

\begin{proposition}[The joint moment generating function, zero start]
\label{prop:mgf}
Let $\lambda=-s$, $z=2t-c^2$, and let $C=\cos\sqrt z$, $S=\sin\sqrt z/\sqrt z$ and $T=\{S+(C-1)/z\}/z^3$, all even in $\sqrt z$ and hence entire functions of $z$.  Then, on the tube,
\begin{equation}
M(s,t)=e^{-s/2}\,D(s,t)^{-1/2},\qquad D(s,t)=e^c\{F_0(z)+F_1(z)\lambda+F_2(z)\lambda^2\},
\label{eq:mgf}
\end{equation}
where, in case \Fk,
\[
F_0=C-cS,\qquad F_1=S,\qquad F_2=0;
\]
in case \Fu,
\begin{gather*}
F_2=\frac{S}{z}+\frac{2(C-1)}{z^2},\qquad F_1=-c^2F_2,\\
F_0=\frac{(2t+c^3)S}{z}-\frac{c^2C}{z}+\frac{4ct(C-1)}{z^2};
\end{gather*}
and in case \Ft,
\begin{align*}
F_2&=-\frac{4(c^2-3c+27)S}{z^2}+\{192t+288(1-c)\}T-\frac{32(c^2-3c+6)C}{z^3}-\frac{16(c^2-3c-3)}{z^3},\\
F_1&=c^4\Big\{\frac{S}{z^2}-24T+\frac{8C}{z^3}+\frac{4}{z^3}\Big\},\\
F_0&=\frac{\{-c^5-8(c^2+3c-3)t\}S}{z^2}-96(1-c)t^2T+\Big\{\frac{c^4}{z^2}-\frac{16c^3t}{z^3}\Big\}C+\frac{8c^2(3-c)t}{z^3}.
\end{align*}
Each $F_m$ is an entire function of $z$, and the square root in (\ref{eq:mgf}) is the branch that is continuous on the tube and positive at real arguments.
\end{proposition}

At $c=0$ in case \Fk, $D=\cosh\mu-(s/\mu)\sinh\mu$ with $\mu=\sqrt{-2t}$, the formula of \cite{Phillips1987}.  The determinants of cases \Fu\ and \Ft\ are rearrangements of Theorems~1 and 2 of \cite{NabeyaSorensen1994}, who write the characteristic function of $xV-U$ for all $x$ as $D_j(2\iu\theta)^{-1/2}e^{\iu\theta}$ with $D_j$ explicit combinations of $\cosh\mu$, $\sinh\mu/\mu$ and powers of $\mu$, $\mu=\sqrt{c^2-2\lambda x}$; since the characteristic function of $xV-U$ for all $x$ is the joint characteristic function, (\ref{eq:mgf}) follows with $\lambda=-s$, $\lambda x=t$ and $D=D_j$, their kernels being twice $U$ and $V$, which does not affect the probability they compute, and the normalization being fixed by the known distribution of $\int_0^1(W-\bar W)^2\dr{r}$.  The form of the proposition exhibits the determinants as polynomials in $\lambda$.  The $F_m$ are evaluated by their Taylor series in $z$ when $|z|<4$, and otherwise from the closed forms written in $\vartheta=(-z)^{1/2}$ and $e^{-2\vartheta}$ with the factor $e^\vartheta$ taken out, so that nothing overflows.

Two properties of $D$ matter.  It is an entire function of $z$, so that the only branch to choose is that of $D^{-1/2}$, which is fixed by continuity from a real point, where $D>0$, and it has no zeros on the tube.  And it is a polynomial in $s$ of degree one (\Fk) or two (\Fu, \Ft), with coefficients that depend on $t$ alone, which is (\ref{eq:U}) in another guise: under the tilt $e^{tV}$ the Gaussian functionals of (\ref{eq:U}) remain Gaussian, with complex covariances, and the moment generating function of a quadratic form of rank $r$ in a Gaussian vector is the inverse square root of a polynomial of degree $r$ in its argument.  The following lemma makes this precise and gives the configuration of the roots along the contours of Section~\ref{sec:inv}.  Let $t_*=\sup\{t:\E e^{tV}<\infty\}$, the first zero of $D(0,\cdot)$ on the positive real axis, and let $(s_2(\tau),s_1(\tau))$, with $s_2(\tau)=-\infty$ in case \Fk, be the interval of convergence of $M(\cdot,\tau)$.

\begin{lemma}[The roots]
\label{lem:roots}
Let $\tau<t_*$ and $\sigma\in(s_2(\tau),s_1(\tau))$.  For every $t$ with $\Rea t=\tau$, $D(\cdot,t)$ is a polynomial in $s$ of exact degree one (\Fk) or two (\Fu, \Ft), with roots $s_1(t),s_2(t)$ continuous in $t$; at $t=\tau$ the roots are real, with $s_2(\tau)<0<s_1(\tau)$ in cases \Fu\ and \Ft\ and $s_1(\tau)>0$ in case \Fk, and $D(\cdot,\tau)>0$ between them; and $\Rea s_2(t)<\sigma<\Rea s_1(t)$, respectively $\sigma<\Rea s_1(t)$, for every $t$ on the line.  Consequently
\begin{equation}
M(s,t)=\begin{cases}e^{-s/2}P(t)\{(s_1-s)(s-s_2)\}^{-1/2}&(\Fu,\Ft),\\[0.5ex]
e^{-s/2}P(t)(s_1-s)^{-1/2}&(\Fk),\end{cases}
\label{eq:Mfactor}
\end{equation}
with $P(t)\ne0$ continuous on the line and $\Rea(s_1-s_2)>0$.
\end{lemma}

In terms of the weights $\kappa_k=1/(2s_k)$, (\ref{eq:Mfactor}) says that, under the tilt $e^{tV}$, $R=U+\frac12$ is distributed as $\kappa_1Z_1^2+\kappa_2Z_2^2$ with $Z_1,Z_2$ independent standard normal, $\kappa_2=0$ in case \Fk; the sign conditions of the lemma are $\Rea\kappa_1>0>\Rea\kappa_2$, respectively $\Rea\kappa_1>0$.  With $\lambda=-s$ the roots are those of $F_0+F_1\lambda+F_2\lambda^2$, so that $s_1,s_2$ and $P(t)=\{-e^cF_2(z)\}^{-1/2}$ (\Fu, \Ft; $e^cF_2<0$ at real points, since $D$ is positive between its roots), $P(t)=\{e^cF_1(z)\}^{-1/2}$ (\Fk), are explicit in the functions of Proposition~\ref{prop:mgf}, and the moment generating function of $V$ is $M_V(t)=M(0,t)=P(t)(-s_1s_2)^{-1/2}$, respectively $P(t)s_1^{-1/2}$.

\subsection{The initial observation}
\label{sec:init}

With a fixed start $\gamma_0$ the Gaussian functionals of (\ref{eq:U}) acquire means proportional to $\gamma_0$, and $R$ becomes a noncentral quadratic form of the same rank; the logarithm of the joint moment generating function gains a term quadratic in $\gamma_0$, which is in closed form.  In cases \Fu\ and \Ft\ the fixed start enters the Fredholm determinants of \citet[Theorems~1 and 2]{NabeyaSorensen1994} through $f=c\gamma_0$ and $f=c^2\gamma_0$, and in case \Fk\ it follows from the Riccati equations of the tilted process (Lemma~\ref{lem:riccati} in Appendix~\ref{app:proofs}).

\begin{proposition}[The joint moment generating function, fixed start]
\label{prop:fixed}
Let $M_0$ and $D$ be the moment generating function and the determinant of the zero start, and let $M_{\gamma_0}$ be the moment generating function of $(U,V)$ with the start $\gamma_0$.  Then, on the tube of $M_0$ and wherever the right side is finite,
\begin{equation}
\log M_{\gamma_0}(s,t)=\log M_0(s,t)+\gamma_0^2\,Q(s,t),
\label{eq:fixedstart}
\end{equation}
where, in cases \Fu\ ($k=1$) and \Ft\ ($k=2$),
\begin{equation}
\begin{gathered}
Q(s,t)=c^{2k}e^c\,\big(t+\tfrac12s^2+cs\big)\,\frac{g(z)}{D(s,t)},\\
g=-\frac{S}{z}-\frac{2(C-1)}{z^2}\ \ (\Fu),\qquad
g=\frac{12C}{z^3}-\frac{(24+4z)(C-1)}{z^4}-\frac{(24-z)S}{z^3}\ \ (\Ft),
\end{gathered}
\label{eq:g}
\end{equation}
both entire in $z$, and in case \Fk, with $F_0$, $F_1$ those of the case,
\begin{equation}
Q(s,t)=\frac{tF_1}{F_0}-\frac{s}{2}+\frac{s}{2F_0^2}\,\frac{1}{1-sF_1/F_0}\qquad(\Fk).
\label{eq:QK}
\end{equation}
As a function of $s$, $Q$ is rational with simple poles at the roots $s_1$, $s_2$ of $D(\cdot,t)$ (at $s_1=F_0/F_1$ in case \Fk), and
\begin{equation}
\begin{gathered}
Q(s,t)-Q(0,t)=\frac{s}{2}\Big\{\frac{\mu_1}{1-2s\kappa_1}-\frac{\mu_2}{1-2s\kappa_2}\Big\},\\
\mu_k=-\frac{2c^{2k}\,(t+\tfrac12s_k^2+cs_k)\,g(z)}{F_2(z)\,(s_1-s_2)\,s_k^2}\qquad(\Fu,\Ft;\ k=1,2),
\end{gathered}
\label{eq:partial}
\end{equation}
and in case \Fk\ $Q(s,t)-Q(0,t)=-\frac s2+\frac s2F_0^{-2}(1-2s\kappa_1)^{-1}$, the constant $-\frac s2$ being the term $-\frac12\gamma_0^2$ of $R=\frac12\{J(1)^2-\gamma_0^2\}$.  Under the tilt $e^{tV}$, $R$ is distributed as $\frac12(X_1^2-X_2^2)$ with $X_1$, $X_2$ independent Gaussian of variances $2\kappa_1$, $-2\kappa_2$ and means $m_k$ with $m_k^2=\gamma_0^2\mu_k$; in case \Fk, $R=\frac12\{J(1)^2-\gamma_0^2\}$ with $J(1)$ of mean $\gamma_0/F_0$ and variance $F_1/F_0$.
\end{proposition}

The partial fractions (\ref{eq:partial}) are the exponent of the moment generating function of a noncentral quadratic form of rank two: under the tilt,
\[
\E\big(e^{sR}\big)=\prod_{k=1}^{2}(1-2s\kappa_k)^{-1/2}\exp\Big\{\frac s2\sum_{k=1}^{2}(-1)^{k-1}\frac{m_k^2}{1-2s\kappa_k}\Big\},
\]
and (\ref{eq:partial}) identifies the squared means, the only way in which the means enter; at real $t$ the $\mu_k$ are nonnegative.  The weights $\kappa_k$ do not depend on the start.  The term $\gamma_0^2Q(0,t)$, the correction to $\log M_V$, is $c^{2k}\gamma_0^2tg(z)/F_0(z)$ in cases \Fu\ and \Ft\ and $\gamma_0^2tF_1/F_0$ in case \Fk.  The parameter $f$ of \cite{NabeyaSorensen1994} collects the initial value and the local drift or trend, $f=|c\gamma_0+d_1|$ and $|c^2\gamma_0+cd_1+d_2|$: the initial value and the drift enter the demeaned or detrended process only through $\gamma_0(e^{cr}-1)$ and $d_1(e^{cr}-1)/c$, and the trend through $d_2(e^{cr}-1-cr)/c^2$, so that (\ref{eq:fixedstart}) also covers a local drift and trend in the data-generating process, with $f^2$ in place of $c^{2k}\gamma_0^2$.

\runin{The stationary start}  With a stationary start the joint moment generating function is the average over $\gamma_0\sim N(0,\sigma_0^2)$, $\sigma_0^2=-1/(2c)$, of (\ref{eq:fixedstart}).  Since $\gamma_0^2\sim\sigma_0^2\chi^2_1$, averaging $e^{\gamma_0^2Q}$ multiplies $M_0$ by $(1-2\sigma_0^2Q)^{-1/2}=(1+Q/c)^{-1/2}$, which amounts to replacing $D$ by
\begin{equation}
D_{\rm S}(s,t)=D(s,t)+c^{2k-1}e^c\,\big(t+\tfrac12s^2+cs\big)\,g(z)\qquad(\Su,\St),
\label{eq:Dstat}
\end{equation}
again a polynomial of degree two in $s$, and in case \Sk\ by
\begin{equation}
D_{\rm S}(s,t)=e^c\Big\{F_0(z)+\frac{S(z)}{2c}\,(2t+s^2)\Big\}\qquad(\Sk),
\label{eq:DstatK}
\end{equation}
of degree two in $s$ with roots $\pm s_2$ and no linear term, as the reversibility of the stationary process requires: $(J(0),J(1))$ is exchangeable, and so is $J(1)^2-J(0)^2$ with its negative, also under the tilt by $\int_0^1J^2\dr{r}$.  Lemma~\ref{lem:roots} holds for $D_{\rm S}$ in all three stationary cases with the configuration of cases \Fu\ and \Ft, and the factorization (\ref{eq:Mfactor}) with $P(t)$ the inverse square root of the leading coefficient; $t_*$ is the first zero of $D_{\rm S}(0,\cdot)$.  Everything that follows applies to the six cases with the appropriate $D$, and to the fixed start with $M_{\gamma_0}$ in place of $M$, where $\varphi_t$ is then the density of the noncentral form of Proposition~\ref{prop:fixed} under the tilt, given in Section~\ref{sec:inv} (Corollary~\ref{cor:fixed}).

\runin{Closed forms free of cancellation}  For $c>0$ the terms of $D(s,t)$ are of the size $e^{2c}$ while the value is of the order of one, and (\ref{eq:mgf}) loses about $2c/\ln10$ digits; at $c=5$ this costs $10^{-12}$, at $c=10$ it costs $10^{-7}$, and at $c=20$ the determinant is lost altogether.  Girsanov's theorem gives forms without this cancellation, which we use for $c>2$.  Under the law of $J$ with a zero start, a functional $\Phi$ of the path has $\E_c\Phi(J)=\E\{\Phi(W)\exp[\frac c2\{W(1)^2-1\}-\frac{c^2}2\int_0^1W(r)^2\dr{r}]\}$.  Let $\bX=W(1)$ in case \Fk, $\bX=(W(1),\int_0^1W\dr{r})'$ in case \Fu\ and $\bX=(W(1),\int_0^1W\dr{r},\int_0^1rW\dr{r})'$ in case \Ft.  Then $V=\int_0^1W(r)^2\dr{r}-\bX'\bGam\bX$ with $\bGam=0$, $\diag(0,1)$ and $\diag(0,\bG^{-1})$, $\bG^{-1}$ the inverse of the Gram matrix of $(1,r)$ on $[0,1]$, with entries $4,-6,-6,12$, and by (\ref{eq:U}) $R=\frac12\bX'\bK\bX$ with $\bK=1$,
\[
\bK=\begin{pmatrix}1&-1\\-1&0\end{pmatrix}\ \ (\Fu),\qquad
\bK=\begin{pmatrix}1&2&-6\\2&-12&12\\-6&12&0\end{pmatrix}\ \ (\Ft).
\]
With $\vartheta=(c^2-2t)^{1/2}$ and $\bQ=\diag(c,0,\ldots,0)-2t\bGam$, for functions $h$ for which the right side converges,
\begin{equation}
\E\big\{e^{tV}h(R)\big\}=e^{-c/2}\,\E\Big\{\exp\Big(-\frac{\vartheta^2}2\int_0^1W(r)^2\dr{r}\Big)\exp\Big(\frac12\bX'\bQ\bX\Big)h\Big(\frac12\bX'\bK\bX\Big)\Big\},
\label{eq:girsanov}
\end{equation}
and under the probability measure with density $(\cosh\vartheta)^{1/2}\exp(-\frac{\vartheta^2}2\int_0^1W^2\dr{r})$ the vector $\bX$ is Gaussian with mean zero and covariance matrix $\bSig$,
\begin{gather*}
\Sigma_{11}=\frac{\tanh\vartheta}\vartheta,\qquad\Sigma_{12}=\frac{1-\sech\vartheta}{\vartheta^2},\qquad\Sigma_{13}=\Sigma_{22}=\frac{\vartheta-\tanh\vartheta}{\vartheta^3},\\
\Sigma_{23}=\frac{\vartheta^2-2+2\sech\vartheta}{2\vartheta^4},\qquad\Sigma_{33}=\frac{\vartheta^3/3-\vartheta+\tanh\vartheta}{\vartheta^5},
\end{gather*}
even in $\vartheta$, with the values $1$, $\frac12$, $\frac13$, $\frac13$, $\frac5{24}$ and $\frac2{15}$ of Brownian motion at $\vartheta=0$.  It follows that $M_V(t)=e^{-c/2}(\cosh\vartheta\,|\bI-\bQ\bSig|)^{-1/2}$ and that, under the tilt, $R$ is a quadratic form in a Gaussian vector with covariance matrix $\bP^{-1}$, $\bP=\bSig^{-1}-\bQ$.  In case \Fk\ this gives
\begin{equation}
M_V(t)=e^{-c/2}\Big(\frac{\vartheta}{\sinh\vartheta}\Big)^{1/2}\eta^{-1/2},\qquad s_1=\eta=\vartheta\coth\vartheta-c,
\label{eq:dkgirs}
\end{equation}
and in cases \Fu\ and \Ft\ the weights $2\kappa_1$, $2\kappa_2$ are the nonzero eigenvalues of $\bK\bP^{-1}$, that is the roots of $\alpha_2\nu^2-\alpha_1\nu+\alpha_0=0$ with $\alpha_2=|\bP|$, $\alpha_1=\tr\{\bK\adj(\bP)\}$ and $\alpha_0=-\tr\{\adj(\bK)\bP\}$ ($\alpha_0=1$ in case \Fu), so that
\begin{equation}
\begin{gathered}
s_1+s_2=-\frac{\alpha_1}{\alpha_0},\qquad s_1-s_2=\frac{(\alpha_1^2+4\alpha_0\alpha_2)^{1/2}}{\alpha_0},\\
\frac{M_V(t)}{(-s_1s_2)^{1/2}}=P(t)=\frac{e^{-c/2}}{(\alpha_0\cosh\vartheta\,|\bSig|)^{1/2}}.
\end{gathered}
\label{eq:dugirs}
\end{equation}
All quantities are computed without cancellation for large $c$: the entries of $\bSig$ from $e^{-2\vartheta}$, the difference $\vartheta-c$ as $-2t/(\vartheta+c)$, and $\eta=(\vartheta-c)+2\vartheta e^{-2\vartheta}/(1-e^{-2\vartheta})$.  The point $t_*$, where $|\bI-\bQ\bSig|$ vanishes, enters $P(t)$ and the product $P(t)\varphi_t$ of Section~\ref{sec:inv} only through $\alpha_2$, polynomially, so these are analytic beyond $t_*$, up to the first zero of $\cosh\vartheta\,|\bSig|$, $\alpha_0$ or $\alpha_1^2+4\alpha_0\alpha_2$ on the real axis, which is at most $(c^2+\pi^2/4)/2$; the Supplementary Material uses this.
\section{Inversion formulas and the representation of the limit}
\label{sec:inv}

Let $(X,Y)$ be a random vector with $Y>0$ almost surely and moment generating function $M(s,t)$, and let $Q_0=(X-c_0Y)/\sqrt Y$ for a real constant $c_0$.  Written out, the distribution function of $Q_0$ is a quadruple integral.  The joint density of $(X,Y)$ is the double inversion $(2\pi\iu)^{-2}\iint M(s,t)e^{-sx-ty}\dr{t}\dr{s}$; the change of variables $(x,y)=(pv+c_0v^2,v^2)$, with the Jacobian $2v^2$, turns $\{Q_0\le q\}$ into $\{p\le q\}$, so that the density of $Q_0$ at $q$ is the joint density on the parabola $x=qv+c_0v^2$, $y=v^2$, integrated over $v$; and the distribution function integrates that density over its argument.  The two outer integrations act on the exponential alone and can be done in closed form.  For $b\in\mathbb{C}$ and $\Rea t>0$ let
\begin{equation}
\Xi_m(b,t)=\int_0^\infty v^me^{-bv-tv^2}\dr{v},
\label{eq:Jm}
\end{equation}
so that, with the principal branch of $\sqrt t$,
\begin{equation}
\Xi_0=\frac12\sqrt{\frac\pi t}\,e^{b^2/(4t)}\erfc\Big(\frac{b}{2\sqrt t}\Big),\quad
\Xi_1=\frac{1-b\Xi_0}{2t},\quad \Xi_2=\frac{\Xi_0-b\Xi_1}{2t}.
\label{eq:Jrec}
\end{equation}
On the parabola $sx+ty=qsv+(t+c_0s)v^2$, so the integral over $v$ of $2v^2e^{-sx-ty}$ is $2\Xi_2(qs,t+c_0s)$; and the integral over the argument of the density, from $q$ to $\infty$ when $\Rea s>0$ and from $-\infty$ to $q$ when $\Rea s<0$, replaces $2v^2e^{-sqv}$ by $\pm2ve^{-sqv}/s$, whose integral over $v$ is $\pm2\Xi_1(qs,t+c_0s)/s$.  What remains is a double inversion with an explicit kernel.

\begin{theorem}[Inversion formulas for $(X-c_0Y)/\sqrt Y$]
\label{th:inv}
Let $M$ be finite in a neighborhood of $(\sigma,\tau)$ with $\tau+c_0\sigma>0$, let $I$ be an open interval, and suppose that
\begin{equation}
\int_{-\infty}^{\infty}\int_{-\infty}^{\infty}|M(\sigma+\iu u,\tau+\iu w)|\,e^{-\delta u^2}\dr{u}\dr{w}<\infty\quad\text{for every }\delta>0,
\label{eq:A}
\end{equation}
and that for every compact $K\subset I$
\begin{equation}
\int_{-\infty}^{\infty}\int_{-\infty}^{\infty}|M(\sigma+\iu u,\tau+\iu w)|\,\sup_{q\in K}|\Xi_2(q(\sigma+\iu u),\tau+c_0\sigma+\iu(w+c_0u))|\dr{u}\dr{w}<\infty.
\label{eq:B}
\end{equation}
Then $Q_0$ has a continuous density on $I$, given for $q\in I$ by
\begin{equation}
f_{Q_0}(q)=\frac{1}{(2\pi\iu)^2}\int_{\sigma-\iu\infty}^{\sigma+\iu\infty}\int_{\tau-\iu\infty}^{\tau+\iu\infty}
M(s,t)\,2\Xi_2(qs,t+c_0s)\dr{t}\dr{s}.
\label{eq:dinv}
\end{equation}
If moreover $\sigma>0$ and (\ref{eq:B}) holds with $\Xi_1(qs,t+c_0s)/s$ in place of $\Xi_2(qs,t+c_0s)$, then for $q\in I$
\begin{equation}
\PP(Q_0>q)=\frac{1}{(2\pi\iu)^2}\int_{\sigma-\iu\infty}^{\sigma+\iu\infty}\int_{\tau-\iu\infty}^{\tau+\iu\infty}
M(s,t)\,\frac{2\Xi_1(qs,t+c_0s)}{s}\dr{t}\dr{s},
\label{eq:cinv}
\end{equation}
and if $\sigma<0$ the same integral equals $-\PP(Q_0\le q)$.
\end{theorem}

The kernels are the two closed-form integrations above, and the sign in the last statement is that of the integral over the argument.  The conditions are mild versions of the absolute integrability of $M$ on the contour, and they are needed: for the limits of Section~\ref{sec:model} the density of $U$ is unbounded (at $U=-\frac12$ in case \Fk, where $U=\{W(1)^2-1\}/2$), so that $|M|$ decays only like $|\Ima s|^{-1/2}$ (\Fk) or $|\Ima s|^{-1}$ (\Fu, \Ft) along the contour; in $t$, $|M|$ decays like $\exp(-\frac12\Rea\mu)$ with $\Rea\mu\sim|\Ima t|^{1/2}$, and both conditions hold, with $I=\mathbb{R}$.  The proof is in Appendix~\ref{app:proofs}.  The interval $I$ need not be the whole line: where the density has an isolated singularity, the theorem is applied on the open intervals on either side of it, and the distribution function is continuous across the point as soon as the singularity is integrable; this is how Theorem~\ref{th:single} treats its one exceptional point.

The double integrals (\ref{eq:dinv}) and (\ref{eq:cinv}) can be evaluated by quadrature, but not efficiently: the kernel $\Xi_2(qs,t)$ is in its Gaussian regime where $\Ima s$ and $\Ima t$ have opposite signs and $|\Ima s|$ is up to a few times $|\Ima t|$, and there it oscillates with the phase $q^2(\Ima s)^2/(4\Ima t)$, which no grid resolves in the far field.  The polynomial structure of Lemma~\ref{lem:roots} does a third of the four integrations in closed form and removes the $s$-integral altogether.  The inverse two-sided Laplace transform of (\ref{eq:Mfactor}) in $s$ at the point $x$ is explicit,
\begin{equation}
\frac{1}{2\pi\iu}\int_{\sigma-\iu\infty}^{\sigma+\iu\infty}M(s,t)e^{-sx}\dr{s}=P(t)\,\varphi_t(X),\qquad X=x+\tfrac12,
\label{eq:phit}
\end{equation}
with
\begin{equation}
\varphi_t(X)=\begin{cases}\dfrac{e^{-s_1X}}{(\pi X)^{1/2}}\,\mathbf{1}\{X>0\}&(\Fk),\\[2ex]
\dfrac1\pi\,e^{-(s_1+s_2)X/2}K_0\Big(\dfrac{(s_1-s_2)|X|}{2}\Big)&(\Fu,\Ft),\end{cases}
\label{eq:varphi}
\end{equation}
the first a scaled $\chi^2_1$ density, the second the density of a difference of two scaled $\chi^2_1$ variables, both continued analytically in $s_1$, $s_2$ and $t$; the convergence of the integrals behind (\ref{eq:phit}) requires $\Rea s_2<\sigma<\Rea s_1$ and $\Rea(s_1-s_2)>0$, which Lemma~\ref{lem:roots} provides.  Up to the normalization, $\varphi_t$ is the density of $R$ under the tilt $e^{tV}$: its total mass is $(-s_1s_2)^{-1/2}=M_V(t)/P(t)$, respectively $s_1^{-1/2}$.  Inserting (\ref{eq:phit}) into (\ref{eq:dinv}) and (\ref{eq:cinv}), with the $v$-integral of (\ref{eq:Jm}) written out, gives the density and the distribution function as integrals over $t$ alone.  Let
\begin{equation}
X(v)=c_0v^2+xv+\tfrac12,
\label{eq:Xv}
\end{equation}
the argument of $\varphi_t$ along the parabola $u=xv+c_0v^2$ in the $(u,v)$-plane, on whose boundary $U=x\sqrt V+c_0V$.

\begin{theorem}[The density and the distribution function of the limit]
\label{th:single}
Let $0<\tau<t_*$.  In each of the six cases, with $D$ as in Section~\ref{sec:mgf}, the limit $\tau$ has a continuous distribution function $\cF(x)$ and a density $f(x)$, given by (\ref{eq:single}) at every $x$ at which the integral converges.  It converges for every $x$, and $f$ is continuous, except in case \Fk\ with $c_0>0$ at $x=-\sqrt{2c_0}$, where it diverges: there the density has a logarithmic singularity, and it is continuous on either side.  The formulas are
\begin{equation}
f(x)=\frac{1}{2\pi\iu}\int_{\tau-\iu\infty}^{\tau+\iu\infty}P(t)\,\Phi^f_x(t)\dr{t},\qquad
\cF(x)=\frac{1}{2\pi\iu}\int_{\tau-\iu\infty}^{\tau+\iu\infty}P(t)\,\Phi_x(t)\,\frac{\dr{t}}{t},
\label{eq:single}
\end{equation}
where
\begin{equation}
\begin{aligned}
\Phi^f_x(t)&=\int_0^\infty2v^2e^{-tv^2}\varphi_t\big(X(v)\big)\dr{v},\\
\Phi_x(t)&=\int_0^\infty2tv\,e^{-tv^2}\Psi_t\big(X(v)\big)\dr{v}
=\Psi_t\big(\tfrac12\big)+\int_0^\infty e^{-tv^2}\,(x+2c_0v)\,\varphi_t\big(X(v)\big)\dr{v},
\end{aligned}
\label{eq:Phidef}
\end{equation}
and $\Psi_t(X)=\int_{-\infty}^X\varphi_t(u)\dr{u}$.
\end{theorem}

The second form of $\Phi_x$ is an integration by parts in $v$; the boundary term at $v=0$ is the mass of $\varphi_t$ on $(-\infty,\frac12)$.  At $c_0=0$ the parabola is the line $u=xv$, and $\Phi_x(t)=\int_{-\infty}^{\infty}\varphi_t(u+\frac12)\,tW_x(u,t)\dr{u}$ with $tW_x(u,t)=\int2tve^{-tv^2}\dr{v}$ over $\{v>0:xv\ge u\}$, which is $1-e^{-tu^2/x^2}$ for $u<0$ and $0$ otherwise when $x<0$, and $e^{-t\max(u/x,0)^2}$ when $x>0$.  The convergence of the $v$-integrals at infinity needs $\Rea(t+c_0s_d)>0$, $s_d$ the root that governs the decay of $\varphi_t$ in the direction in which $X(v)$ moves ($s_1$ for $X\to+\infty$, $s_2$ for $X\to-\infty$); this holds on the whole line, because $\Rea s_1>0>\Rea s_2$ there by Lemma~\ref{lem:roots}, and $c_0$ and $s_d$ have the same sign whenever $X\to\pm\infty$ requires it.  Hypothesis (\ref{eq:A}) of Theorem~\ref{th:inv} holds on the tube for $\sigma$ with $\tau+c_0\sigma>0$; hypothesis (\ref{eq:B}) fails at two points, and the proof in Appendix~\ref{app:proofs} goes through the smoothing argument of Theorem~\ref{th:inv} instead.

The transforms (\ref{eq:Phidef}) are special functions of $t$, $s_1(t)$ and $s_2(t)$.  For case \Fk\ let
\begin{equation}
\begin{aligned}
\cP_k(\beta,\gamma)&=\int_0^\infty X^{k-1/2}e^{-\beta X^2-\gamma X}\dr{X}\\
&\qquad=\Gamma(k+\tfrac12)(2\beta)^{-(2k+1)/4}e^{\gamma^2/(8\beta)}D_{-k-1/2}\{\gamma(2\beta)^{-1/2}\},\\
\cQ_k(\alpha,\beta)&=\int_0^\infty u^k(1+u)^{-1/2}e^{-\alpha u^2-\beta u}\dr{u},
\end{aligned}
\label{eq:PQ}
\end{equation}
for $\Rea\beta>0$ and $\Rea\alpha>0$, the first a parabolic cylinder function, which is a Bessel function of order $\frac14$ or $\frac34$, $D_{-1/2}(x)=(x/2\pi)^{1/2}K_{1/4}(x^2/4)$ and $D_{1/2}(x)=\frac12x(x/2\pi)^{1/2}\{K_{1/4}(x^2/4)+K_{3/4}(x^2/4)\}$, with $D_{-3/2}$ and $D_{-5/2}$ from $D_{\nu+1}(x)=xD_\nu(x)-\nu D_{\nu-1}(x)$; the second an incomplete one, with $\cQ_0(0,\beta)=e^\beta(\pi/\beta)^{1/2}\erfc(\beta^{1/2})$ and $\cQ_k=(-\partial_\beta)^k\cQ_0$.  For the other cases let
\begin{equation}
\cG^{\pm}_k(a;t)=\int_{\pm(X-1/2)>0}(X-\tfrac12)^ke^{-a(X-1/2)^2}\varphi_t(X)\dr{X},
\label{eq:G}
\end{equation}
the Gaussian transforms of $\varphi_t$ over the two half-lines, with $\cG^+_0(0;t)+\cG^-_0(0;t)=(-s_1s_2)^{-1/2}$, the total mass of $\varphi_t$.

\begin{corollary}[Closed forms at $c_0=0$]
\label{cor:closed}
Let $c_0=0$ and $y=1/(4x^2)$.  In case \Fk,
\begin{align}
\Phi_x(t)&=\frac{\erf\{(s_1/2)^{1/2}\}}{s_1^{1/2}}+\frac{e^{-s_1/2}}{(2\pi)^{1/2}}\,\cQ_0\Big(\frac{t}{4x^2},\frac{s_1}{2}\Big)-\frac{\mathbf{1}\{x<0\}}{\pi^{1/2}}\,e^{-ty}\,\cP_0\Big(\frac{t}{x^2},s_1-\frac{t}{x^2}\Big),\label{eq:PhiK}\\
\Phi^f_x(t)&=\frac{e^{-s_1/2}}{2x^3(2\pi)^{1/2}}\,\cQ_2\Big(\frac{t}{4x^2},\frac{s_1}{2}\Big)+\frac{2\,\mathbf{1}\{x<0\}}{|x|^3\pi^{1/2}}\,e^{-ty}\Big(\cP_2-\cP_1+\frac{\cP_0}{4}\Big)\Big(\frac{t}{x^2},s_1-\frac{t}{x^2}\Big).\label{eq:PhifK}
\end{align}
In cases \Fu\ and \Ft, and in the three stationary cases,
\begin{equation}
\begin{gathered}
\Phi_x(t)=\begin{cases}(-s_1s_2)^{-1/2}-\cG^+_0(0;t)+\cG^+_0(t/x^2;t)&(x>0),\\
\cG^-_0(0;t)-\cG^-_0(t/x^2;t)&(x<0),\end{cases}\qquad
\Phi^f_x(t)=\frac{2\,\cG^{\pm}_2(t/x^2;t)}{|x|^3},
\end{gathered}
\label{eq:PhiUT}
\end{equation}
with $\cG^+$ for $x>0$ and $\cG^-$ for $x<0$.
\end{corollary}

The structure of (\ref{eq:PhiK}) is that of the general case.  The terms without the factor $e^{-ty}$ are analytic in $t$ and vary on the scale of $t$ itself; the terms with it, present for $x<0$ only, are the contributions of a neighborhood of the support edge $X=0$ of $\varphi_t$, where the kernel $e^{-tu^2/x^2}$ takes the value $e^{-ty}$, and they oscillate in $\Ima t$ with the period $2\pi/y=8\pi x^2$.  In the other cases the edge is the logarithmic singularity of $K_0$ at $X=0$, and the same two parts are present inside $\cG^-_0(t/x^2;t)$ and $\cG^-_2(t/x^2;t)$.  With $c_0\ne0$ the edge $X=0$ is met at the positive roots of $X(v)$, of which there are none, one or two: for $c_0<0$ exactly one, for every $x$, so that the oscillating part is present also for $x>0$; for $c_0>0$ two when $x<-\sqrt{2c_0}$, between which $X<0$ and $\varphi_t$ vanishes in case \Fk, and none otherwise.  Each root $v_j$ carries the factor $e^{-tv_j^2}$, oscillating in $\Ima t$ with the period $2\pi/v_j^2$.  In case \Fk\ with $c_0>0$ the two roots merge at $x=-\sqrt{2c_0}$, where the density has a logarithmic singularity: the $t$-ratio centered at an explosive null value has an infinite density at the point $-\sqrt{2c_0}$, the largest value that $\tau=-(\frac12+c_0V)/\sqrt V$ takes on the event $J(1)=0$, where $R=0$.  The separation of the two parts, and the evaluation of all the transforms to twelve digits, is the subject of the Supplementary Material.

\runin{The initial observation}  With a fixed start $\gamma_0$, (\ref{eq:phit}) holds with $M_{\gamma_0}$ in place of $M$, and $\varphi_t$ is the density of the noncentral form of Proposition~\ref{prop:fixed} under the tilt, normalized as before.  In case \Fk, $R+\frac12\gamma_0^2=\frac12J(1)^2$ with $J(1)\sim N(m,1/s_1)$, $m=\gamma_0/F_0$, so that, with $\tilde X=X+\frac12\gamma_0^2$,
\begin{equation}
\varphi_t(X)=e^{\gamma_0^2tF_1/F_0}\,e^{-s_1m^2/2}\,\frac{e^{-s_1\tilde X}\cosh\{ms_1(2\tilde X)^{1/2}\}}{(\pi\tilde X)^{1/2}}\,\mathbf{1}\{\tilde X>0\},
\label{eq:phinc}
\end{equation}
the support edge having moved to $\tilde X=0$, and $\Psi_t(\frac12)=e^{\gamma_0^2tF_1/F_0}s_1^{-1/2}\bar F\{(1+\gamma_0^2)^{1/2};m,1/s_1\}$ with $\bar F(r;m,w)=\frac12\{\erf[(r-m)/(2w)^{1/2}]+\erf[(r+m)/(2w)^{1/2}]\}$, the probability that $X^2\le r^2$ for $X\sim N(m,w)$.  In cases \Fu\ and \Ft, $R=\frac12(X_1^2-X_2^2)$ with $X_1\sim N(m_1,1/s_1)$ and $X_2\sim N(m_2,-1/s_2)$ independent, $m_k^2=\gamma_0^2\mu_k$, and integrating out $X_2$ gives
\begin{equation}
\begin{gathered}
\varphi_t(X)=\frac2\pi\,e^{\gamma_0^2Q(0,t)}\,e^{(s_2m_2^2-s_1m_1^2)/2}\,e^{-s_dX}\,G(\pm X;s_1-s_2,b_1,b_2),\\
G(Y;a,b_1,b_2)=\int_0^\infty e^{-al^2/2}\cosh(b_1l)\,\frac{\cosh\{b_2(2Y+l^2)^{1/2}\}}{(2Y+l^2)^{1/2}}\dr{l},
\end{gathered}
\label{eq:phinc2}
\end{equation}
with the upper sign, $s_d=s_1$, $b_1=m_2s_2$ and $b_2=m_1s_1$ for $X>0$, and the lower sign, $s_d=s_2$, $b_1=m_1s_1$ and $b_2=m_2s_2$ for $X<0$; $G(Y;a,0,0)=\frac12e^{aY/2}K_0(aY/2)$ gives back (\ref{eq:varphi}), the sign of the means is immaterial, and $G$ is analytic in $Y$ on $\Rea Y>0$, which is where the quadrature of the Supplementary Material evaluates it.  The kink of $\varphi_t$ stays at $X=0$.  What changes is that $\varphi_t$ is a Gaussian in $(2\tilde X)^{1/2}$ centered at the mean, so that its mass sits at a distance from the edge, by the factor $e^\lambda$ with $\lambda=m^2s_d/2$ ($12$ at $c=-5$ and $\gamma_0=1$ in case \Fu); the quadrature of the Supplementary Material allows for this.

\begin{corollary}[The fixed start]
\label{cor:fixed}
With a fixed start $\gamma_0$ in cases \Fk, \Fu\ and \Ft, Theorem~\ref{th:single} holds with $M_{\gamma_0}$ in place of $M$ and with $\varphi_t$ the density (\ref{eq:phinc}) or (\ref{eq:phinc2}), except that in case \Fk\ with $c_0>0$ the density has its logarithmic singularity at
\begin{equation*}
x=-\{2c_0(1+\gamma_0^2)\}^{1/2},
\end{equation*}
the largest value of $\tau$ on the event $J(1)=0$.
\end{corollary}

\section{Moments}
\label{sec:mom}

The moments of the limits need no inversion.  The gamma integral
\[
V^{-p/2}=\frac{1}{\Gamma(p/2)}\int_0^\infty s^{p/2-1}e^{-sV}\dr{s}
\]
turns them into one-dimensional integrals of the moment generating function at real arguments.

\begin{proposition}[Moments of the limit]
\label{prop:mom}
Let $\tau=(U-c_0V)/\sqrt V$ and let $p\ge1$ be an integer.  Then
\begin{equation}
\E\big(\tau^p\big)=\frac{1}{\Gamma(p/2)}\int_0^\infty s^{p/2-1}\,m_p(s)\dr{s},\qquad
m_p(s)=\frac{\partial^pM(u,-s-c_0u)}{\partial u^p}\Big|_{u=0},
\label{eq:mom}
\end{equation}
with $M$ the moment generating function of the case: (\ref{eq:mgf}) with a zero start, (\ref{eq:fixedstart}) with a fixed start, and (\ref{eq:mgf}) with $D_{\rm S}$ in place of $D$ with a stationary start.
\end{proposition}

The moments exist for every $p$: by the Cauchy--Schwarz inequality
\[
\E\big(|U-c_0V|^pe^{-sV}\big)\le\big\{\E(U-c_0V)^{2p}\,M(0,-2s)\big\}^{1/2},
\]
and $M(0,-2s)=D(0,-4s-c^2)^{-1/2}$ decays like $e^{-s^{1/2}}$ up to a power of $s$, so the integral of $s^{p/2-1}\E(|U-c_0V|^pe^{-sV})$ converges at both ends; Fubini's theorem gives (\ref{eq:mom}) from the gamma integral, and differentiation under the expectation, which the moment generating function permits near $u=0$, gives the inner derivative.  The moment generating function is analytic in $u$ in the disc $|u|<\min_k|s_k(-s)|$ around $0$, by Lemma~\ref{lem:roots} applied on the real line $t=-s$, and the derivatives are computed as Cauchy integrals on a circle inside it, the integral over $s$ by the trapezoidal rule in $\log s$; for $c>0$ the Girsanov forms (\ref{eq:girsanov})--(\ref{eq:dugirs}) replace (\ref{eq:mgf}).  The quadrature is in the Supplementary Material; a set of four moments takes a fraction of a second.

\section{Numerical results}
\label{sec:results}

This section sets the limits of Theorem~\ref{th:single} and the moments of Proposition~\ref{prop:mom} beside the critical values, the exact results and the simulations in the literature.  The accuracy and the cost of the programs, and the checks against known values, are in the Supplementary Material.

\subsection{Quantiles and critical values}
\label{sec:quant}

Table~\ref{tab:quant} gives the quantiles of the three limiting distributions at the unit root and under the local alternative $c=-5$, with a zero start, and the asymptotic critical values of \cite{MacKinnon2010}, which agree to within $1.3\times10^{-4}$, consistent with their standard errors.  The Richardson extrapolation of the exact finite-sample quantiles of \cite{BaoBrodaKan2026} from $n=500$, $1000$ and $2000$ is within $10^{-5}$ of these in 22 of the 27 cases at $c=0$, and off by up to $6\times10^{-5}$ in the upper tails, mainly of $\hat\tau$, where the finite-sample quantiles approach their limit more slowly.

\begin{table}[tbp]
\caption{Quantiles of the limiting Dickey--Fuller $t$-ratios}
\label{tab:quant}
\begin{center}
\scriptsize
\setlength{\tabcolsep}{1.8pt}
\begin{tabular}{llrrrrrrrrrr}
\toprule
 & $c$ & 1\% & 2.5\% & 5\% & 10\% & 25\% & 50\% & 90\% & 95\% & 97.5\% & 99\%\\
\midrule
$\hat\tau$ & $0$ & $-2.5658$ & $-2.2274$ & $-1.9409$ & $-1.6167$ & $-1.0913$ & $-0.5001$ & $0.8875$ & $1.2831$ & $1.6230$ & $2.0145$\\
\quad M10 & & $-2.5657$ & & $-1.9410$ & $-1.6168$ & & & & & & \\
 & $-5$ & $-3.2587$ & $-2.9735$ & $-2.7368$ & $-2.4751$ & $-2.0674$ & $-1.6623$ & $-1.0126$ & $-0.8321$ & $-0.6579$ & $-0.4310$\\
$\hat\tau_\mu$ & $0$ & $-3.4304$ & $-3.1221$ & $-2.8616$ & $-2.5668$ & $-2.0874$ & $-1.5664$ & $-0.4403$ & $-0.0785$ & $0.2380$ & $0.6067$\\
\quad M10 & & $-3.4303$ & & $-2.8615$ & $-2.5668$ & & & & & & \\
 & $-5$ & $-3.7216$ & $-3.4345$ & $-3.1943$ & $-2.9258$ & $-2.4999$ & $-2.0632$ & $-1.3233$ & $-1.1126$ & $-0.9143$ & $-0.6634$\\
$\hat\tau_\tau$ & $0$ & $-3.9587$ & $-3.6614$ & $-3.4105$ & $-3.1270$ & $-2.6685$ & $-2.1806$ & $-1.2464$ & $-0.9400$ & $-0.6597$ & $-0.3251$\\
\quad M10 & & $-3.9588$ & & $-3.4105$ & $-3.1271$ & & & & & & \\
 & $-5$ & $-4.1316$ & $-3.8452$ & $-3.6044$ & $-3.3338$ & $-2.9002$ & $-2.4477$ & $-1.6594$ & $-1.4348$ & $-1.2267$ & $-0.9653$\\
\bottomrule
\end{tabular}
\end{center}
{\small Zero start, no drift, $c_0=0$.  Quantiles of the limits of $\hat\tau$, $\hat\tau_\mu$ and $\hat\tau_\tau$ under $\phi=1+c/n$ from Theorem~\ref{th:single}; M10: the asymptotic critical values of \cite{MacKinnon2010}.  To five digits, at the probabilities $0.01$, $0.025$, $0.05$, $0.10$, $0.50$, $0.90$, $0.95$, $0.975$, $0.99$ and $c=0$:
$-2.56580$, $-2.22741$, $-1.94087$, $-1.61670$, $-0.50014$, $0.88754$, $1.28314$, $1.62295$, $2.01451$ ($\hat\tau$),
$-3.43035$, $-3.12209$, $-2.86157$, $-2.56679$, $-1.56643$, $-0.44030$, $-0.07847$, $0.23801$, $0.60667$ ($\hat\tau_\mu$),
$-3.95868$, $-3.66140$, $-3.41045$, $-3.12702$, $-2.18060$, $-1.24636$, $-0.94004$, $-0.65970$, $-0.32511$ ($\hat\tau_\tau$).}
\end{table}

\subsection{Comparison with \cite{Abadir1995} and \cite{Larsson1995}}
\label{sec:larsson}

At $c=c_0=0$ in case \Fk, Theorem~\ref{th:single} can be compared with \cite{Abadir1995}.  His starting point, the joint density after the inversion in the argument of $U$ (his (2.13)), is (\ref{eq:phit})--(\ref{eq:varphi}) at $c=0$.  He then carries out the inversion in $t$ analytically: for $x<0$ by a binomial expansion, which gives a convergent triple series of incomplete gamma functions (his (2.11)), and for $x\ge0$ by changing the order of integration, which leaves one integral of a parabolic cylinder function (his (2.17) and (2.18)).  Evaluated in 30-digit arithmetic, his distribution function and density agree with (\ref{eq:single}) to $6\times10^{-16}$ at sixteen points of $[-4,3]$, and his mode of the density, $-0.724$, with ours, $-0.7235$.  One statement of his does not survive the comparison: his (3.6), repeated from \citet[p.~316]{Abadir1992}, puts the density of $\hat\tau$ at zero at $(\pi e)^{-1/2}=0.3422$, on the argument that the densities of $R/\sqrt S$ and of $R/S$ at zero both equal that of $R$ at zero; for a ratio, $f_{R/\sqrt S}(0)=\E(\sqrt S\mid R=0)f_R(0)$, and the value is $f(0)=0.325567$, which his own (2.18) gives as well at $x=0.001$.

Table~\ref{tab:larsson} sets the limit against the values that \cite{Larsson1995} tabulated for case \Fk, for the Dickey--Fuller statistic under the explosive local alternative $\phi=e^{\beta/n}$ and for the $t$-ratio centered at the true value, $c=c_0=\beta$.  His simulated values, from $10^5$ series of length $400$, agree with the limit to within $0.006$ and mostly within their sampling error; the two largest differences, $0.84$ against $0.846$ at $\beta=1$, $z=2$ in panel A and $0.071$ against $0.0747$ at $\beta=2$, $z=-2$ in panel B, are of four to five standard errors (in the latter the exact finite-sample probability at $n=400$ is $0.0756$ \citep{BaoBrodaKan2026}, so the finite sample does not account for it).  His approximation (3.4) to his exact upper-tail formula is accurate only far in the upper tail: at $\beta=2$ it gives $0.05$ and $0.14$ at $z=0.5$ and $1$, where the limit is $0.277$ and $0.343$.  His saddlepoint tail approximations are close in the lower tail at $\beta=1$ ($0.029$ for $0.026$ at $z=-2$) but not at $\beta=2$ ($0.003$ for $0.015$), and in panel B they miss the upper tail ($0.59$ and $0.92$ for $0.961$ and $0.987$ at $\beta=2$).

\begin{table}[tbp]
\caption{The limiting distribution of $\hat\tau$ under local alternatives: \cite{Larsson1995} and the exact values}
\label{tab:larsson}
\begin{center}
\footnotesize
\setlength{\tabcolsep}{4pt}
\begin{tabular}{rrrrrcrrrr}
\toprule
\multicolumn{10}{l}{\emph{A. Dickey--Fuller test, $\PP(\hat\tau\le z)$ under $\phi=e^{\beta/n}$ ($c=\beta$, $c_0=0$; his Table B.3)}}\\[0.3ex]
 & \multicolumn{4}{c}{$\beta=1$} & & \multicolumn{4}{c}{$\beta=2$}\\
\cmidrule(lr){2-5}\cmidrule(lr){7-10}
$z$ & Sim. & (3.4) & Tail & Exact & & Sim. & (3.4) & Tail & Exact\\
\midrule
$-2$   & 0.026 & & 0.029 & 0.026405 & & 0.015 & & 0.003 & 0.015333\\
$-1.5$ & 0.076 & & 0.096 & 0.075778 & & 0.042 & & 0.019 & 0.043064\\
$-1$   & 0.17  & & 0.30  & 0.171927 & & 0.093 & &       & 0.094033\\
$0$    & 0.42  & &       & 0.424176 & & 0.21  & &       & 0.215288\\
$0.5$  & 0.55  & 0.47 & & 0.545650 & & 0.28  & 0.05 & & 0.277173\\
$1$    & 0.66  & 0.61 & & 0.661230 & & 0.34  & 0.14 & & 0.342599\\
$2$    & 0.84  & 0.83 & & 0.845905 & & 0.48  & 0.36 & & 0.479265\\
$3$    & 0.948 & 0.945 & & 0.947926 & & 0.61 & 0.54 & & 0.611544\\
$5$    & 0.998 & 0.998 & & 0.997702 & & 0.82 & 0.80 & & 0.818787\\
$8$    & 1     & 1     & & 0.999998 & & 0.964 & 0.960 & & 0.962540\\
\midrule
\multicolumn{10}{l}{\emph{B. $t$-ratio centered at the true value, $\PP(\hat\tau_0\le z)$ with $c=c_0=\beta$ (his Table B.2)}}\\[0.3ex]
 & \multicolumn{4}{c}{$\beta=-2$} & & \multicolumn{4}{c}{$\beta=2$}\\
\cmidrule(lr){2-5}\cmidrule(lr){7-10}
$z$ & Sim. & & Tail & Exact & & Sim. & & Tail & Exact\\
\midrule
$-2$   & 0.039 & & 0.040 & 0.039149 & & 0.071 & & 0.081 & 0.074729\\
$-1.5$ & 0.11  & & 0.10  & 0.109743 & & 0.19  & & 0.33  & 0.191272\\
$-1$   & 0.24  & &       & 0.244738 & & 0.32  & &       & 0.322495\\
$0$    & 0.65  & &       & 0.647192 & & 0.65  & &       & 0.650436\\
$1$    & 0.92  & &       & 0.917900 & & 0.90  & &       & 0.901895\\
$1.5$  & 0.97  & & 0.61  & 0.969410 & & 0.96  & & 0.59  & 0.960805\\
$2$    & 0.99  & & 0.94  & 0.990543 & & 0.99  & & 0.92  & 0.987266\\
$3$    & 0.999 & & 0.998 & 0.999515 & & 0.999 & & 0.998 & 0.999303\\
\bottomrule
\end{tabular}
\end{center}
{\small Case \Fk, zero start.  Sim.: \cite{Larsson1995}, Tables B.2 and B.3, column~1, from $10^5$ simulated series of length 400 (standard error at most $1.5\times10^{-3}$); (3.4): his approximation (3.4) to his exact upper-tail formula (3.3); Tail: his saddlepoint tail approximations (Theorem~4.1); blank where he gives none.  Exact: Theorem~\ref{th:single}.}
\end{table}

\subsection{Explosive alternatives and the right-tailed test}
\label{sec:bubble}

Tests for speculative bubbles use the other tail of the same statistic.  \cite{DibaGrossman1988} and, in the form now standard, \cite{PhillipsWuYu2011} and \cite{PhillipsShiYu2015} test the unit root against the explosive alternative $\phi>1$ by rejecting for large values of $\hat\tau_\mu$, over the full sample and over recursive or rolling windows.  For the test on a fixed window the null limit is that of Table~\ref{tab:quant}, so its critical values are the upper quantiles there, $-0.0785$ and $0.6067$ at the $5\%$ and $1\%$ levels for $\hat\tau_\mu$, which are the values that \cite{PhillipsWuYu2011} obtain by simulation to two decimals, and its limiting power against the explosive local alternatives $\phi=1+c/n$, $c>0$, is the upper tail of the distributions of Section~\ref{sec:inv}, which cover $c>0$ as they cover $c<0$.  Table~\ref{tab:bubble} gives it.  The power rises steeply between $c=1$ and $c=3$, from a fifth to four fifths without deterministic terms, and the test with a trend is blind to the alternative until $c$ exceeds about $1$: at $c=0.5$ and $1$ its power is below its size, the explosive drift being absorbed by the trend.  The sup statistics of \cite{PhillipsWuYu2011}, which take the maximum over windows, have a different limit, a supremum of the functionals (\ref{eq:limit}) over subintervals, which the present method does not reach.

\begin{table}[tbp]
\caption{Limiting power of the right-tailed Dickey--Fuller tests against explosive alternatives}
\label{tab:bubble}
\begin{center}
\small
\setlength{\tabcolsep}{5pt}
\begin{tabular}{llrrrrrr}
\toprule
 & Level & Critical value & \multicolumn{5}{c}{$c$}\\
\cmidrule(lr){4-8}
 & & & $0.5$ & $1$ & $2$ & $3$ & $5$\\
\midrule
$\hat\tau$ & 5\% & $1.2831$ & 0.134 & 0.279 & 0.619 & 0.835 & 0.973\\
 & 1\% & $2.0145$ & 0.047 & 0.152 & 0.519 & 0.792 & 0.968\\[0.4ex]
$\hat\tau_\mu$ & 5\% & $-0.0785$ & 0.087 & 0.170 & 0.569 & 0.844 & 0.978\\
 & 1\% & $0.6067$ & 0.022 & 0.060 & 0.423 & 0.792 & 0.972\\[0.4ex]
$\hat\tau_\tau$ & 5\% & $-0.9400$ & 0.048 & 0.048 & 0.225 & 0.707 & 0.970\\
 & 1\% & $-0.3251$ & 0.010 & 0.010 & 0.109 & 0.616 & 0.961\\
\bottomrule
\end{tabular}
\end{center}
{\small $\PP(\tau>x)$ under $\phi=1+c/n$, zero start, $c_0=0$, with $x$ the upper $5\%$ and $1\%$ quantiles of the null limit (Table~\ref{tab:quant}), from Theorem~\ref{th:single}.}
\end{table}

\subsection{Local power and the initial observation}
\label{sec:power}

Table~\ref{tab:init} gives the limiting power of the $5\%$ tests, with the limiting critical values of Table~\ref{tab:quant}, against $\phi=1+c/n$ with a fixed start of $\xi$ stationary standard deviations ($\gamma_0=\xi/\sqrt{-2c}$) and with a stationary start; for $\hat\tau$, whose null distribution itself depends on a start of this order, the critical value is held at its zero-start value.  Four features stand out.
\begin{itemize}
\item Power rises with the initial deviation, and most without deterministic terms: at $c=-5$ it rises from $0.318$ at $\xi=0$ to $0.720$ at $\xi=3$ for $\hat\tau$, from $0.117$ to $0.208$ for $\hat\tau_\mu$, and only from $0.083$ to $0.090$ for $\hat\tau_\tau$.  This agrees with the power curves of \cite{MullerElliott2003} and the simulations of \cite{HarveyLeybourneTaylor2009}, in which the power of the least squares Dickey--Fuller tests grows with the magnitude of the initial deviation, unlike that of the GLS-detrended tests.
\item Close to the unit root the effect reverses for $\hat\tau_\tau$: at $c=-2.5$ its power falls from $0.0595$ at $\xi=0$ to $0.0570$ at $\xi=3$, and from $0.0595$ to $0.0592$ with a stationary start.
\item The stationary start behaves almost exactly like the fixed start $\xi=1$, its root mean square: at $c=-5$ the powers are $0.362$ and $0.361$ for $\hat\tau$ and $0.125$ and $0.124$ for $\hat\tau_\mu$.
\item The stationary-start powers reproduce the simulations of \cite{Elliott1999}.  His large-sample powers of the Dickey--Fuller $t$-tests with a stationary start (his Table~2, from $T=1000$ and $20{,}000$ replications) are $0.13$, $0.33$, $0.62$ and $0.86$ with an intercept and $0.08$, $0.19$, $0.39$ and $0.63$ with a trend, at $c=-5$, $-10$, $-15$ and $-20$.  The limits are $0.125$, $0.325$, $0.618$ and $0.859$, and $0.084$, $0.188$, $0.379$ and $0.621$.
\end{itemize}
The last rows of each panel are the exact powers at $n=100$ with a stationary start from \cite{BaoBrodaKan2026}, which use the exact finite-sample critical values.  The limit is an excellent guide to the size-adjusted power at $n=100$ near the unit root, within $0.005$ for $c\ge-5$, and it understates power further away, by up to $0.045$ for $\hat\tau_\tau$ at $c=-20$; for $\hat\tau_\tau$ at $c=-15$ and $-20$ the values of \cite{Elliott1999} lie between the limit and the exact power at $n=100$ ($0.404$ and $0.664$), as a sample of $T=1000$ would.

\begin{table}[tbp]
\caption{Limiting power of the 5\% Dickey--Fuller tests and the initial observation}
\label{tab:init}
\begin{center}
\small
\setlength{\tabcolsep}{5pt}
\begin{tabular}{llrrrrrr}
\toprule
 & Start & \multicolumn{6}{c}{$c$}\\
\cmidrule(lr){3-8}
 & & $-30$ & $-20$ & $-15$ & $-10$ & $-5$ & $-2.5$\\
\midrule
$\hat\tau$ & $\xi=0$ & 1.000 & 0.998 & 0.967 & 0.758 & 0.318 & 0.144 \\
 & $\xi=1$ & 1.000 & 0.999 & 0.974 & 0.792 & 0.361 & 0.173 \\
 & $\xi=2$ & 1.000 & 0.999 & 0.988 & 0.877 & 0.499 & 0.277 \\
 & $\xi=3$ & 1.000 & 1.000 & 0.997 & 0.958 & 0.720 & 0.499 \\
 & stationary & 1.000 & 0.998 & 0.973 & 0.789 & 0.362 & 0.176 \\
 & stationary, $n=100$ & 1.000 & 0.999 & 0.976 & 0.802 & 0.369 & 0.178 \\[0.4ex]
$\hat\tau_\mu$ & $\xi=0$ & 0.995 & 0.844 & 0.593 & 0.304 & 0.117 & 0.072 \\
 & $\xi=1$ & 0.996 & 0.860 & 0.619 & 0.325 & 0.124 & 0.074 \\
 & $\xi=2$ & 0.998 & 0.901 & 0.691 & 0.391 & 0.152 & 0.081 \\
 & $\xi=3$ & 0.999 & 0.948 & 0.795 & 0.507 & 0.208 & 0.098 \\
 & stationary & 0.996 & 0.859 & 0.618 & 0.325 & 0.125 & 0.074 \\
 & stationary, $n=100$ & 0.997 & 0.883 & 0.646 & 0.337 & 0.126 & 0.074 \\[0.4ex]
$\hat\tau_\tau$ & $\xi=0$ & 0.940 & 0.606 & 0.366 & 0.182 & 0.083 & 0.059 \\
 & $\xi=1$ & 0.945 & 0.621 & 0.379 & 0.188 & 0.084 & 0.059 \\
 & $\xi=2$ & 0.959 & 0.666 & 0.417 & 0.206 & 0.086 & 0.058 \\
 & $\xi=3$ & 0.975 & 0.734 & 0.483 & 0.241 & 0.090 & 0.057 \\
 & stationary & 0.945 & 0.621 & 0.379 & 0.188 & 0.084 & 0.059 \\
 & stationary, $n=100$ & 0.965 & 0.664 & 0.404 & 0.195 & 0.084 & 0.059 \\
\bottomrule
\end{tabular}
\end{center}
{\small $\PP(\tau\le x_{0.05})$ under $\phi=1+c/n$, no drift, $c_0=0$, with $x_{0.05}$ the limiting 5\% critical value at the unit root with a zero start ($-1.94087$, $-2.86157$, $-3.41045$).  $\xi$: fixed start $y_0=\sigma\sqrt n\,\xi/\sqrt{-2c}$, from Theorem~\ref{th:single}, with (\ref{eq:phinc}) and (\ref{eq:phinc2}) for $\xi>0$; stationary: $y_0$ from the stationary distribution, with $D_{\rm S}$ of Section~\ref{sec:init}; stationary, $n=100$: the exact finite-sample power of \cite{BaoBrodaKan2026} with the exact critical values.}
\end{table}

\subsection{Moments}
\label{sec:momnum}

Table~\ref{tab:mom} gives the mean and standard deviation of the limit of the Dickey--Fuller statistic under local alternatives, with a zero and with a stationary start, from Proposition~\ref{prop:mom}.  At the unit root the first six moments of the three limits agree with the eight-digit values of \cite{Nabeya1999}, computed from the formula of \cite{Sawa1978} with Simpson's rule, to within one unit in their last digit; the second moment of $\hat\tau_\mu$, $3.05599535$, is also the expectation $3.0560$ of the limiting likelihood ratio statistic in \cite{Larsson1997}.  Under local alternatives the first four moments agree with Gauss--Legendre integration of the densities of Theorem~\ref{th:single} to $2\times10^{-11}$ relative ($\hat\tau_\mu$ at $c=-5$, $\hat\tau_\tau$ at $c=-10$, $\hat\tau_\tau$ at $c=c_0=-5$, and with a stationary start), and on the explosive side to $1.2\times10^{-10}$ at $c=c_0=10$; on that side the means of the Dickey--Fuller statistic grow fast, $59.0696374465$ and $45.7021278187$ in cases \Fk\ and \Fu\ at $c=5$.  As $c$ decreases the distribution moves to the left and becomes narrower, and the shift is largest without deterministic terms: from $c=0$ to $c=-5$ the mean falls by $1.28$ for $\hat\tau$, by $0.56$ for $\hat\tau_\mu$ and by $0.29$ for $\hat\tau_\tau$, which is the ordering of the local powers in Table~\ref{tab:init}.  A stationary start moves the mean further to the left for $\hat\tau$ at every $c$, and for $\hat\tau_\mu$ and $\hat\tau_\tau$ from $c=-5$ and $c=-10$ on; closer to the unit root it moves them slightly to the right, the counterpart of the small loss of power there.  It changes the standard deviation by at most $0.016$.

\begin{table}[tbp]
\caption{Mean and standard deviation of the limit of the Dickey--Fuller $t$-ratio under local alternatives}
\label{tab:mom}
\begin{center}
\small
\begin{tabular}{llrrrrr}
\toprule
 & Start & \multicolumn{5}{c}{$c$}\\
\cmidrule(lr){3-7}
 & & $0$ & $-2.5$ & $-5$ & $-10$ & $-20$\\
\midrule
$\hat\tau$ & zero & $-0.4231$ & $-1.2636$ & $-1.7066$ & $-2.3352$ & $-3.2366$\\
 & & (0.9811) & (0.6548) & (0.5867) & (0.5461) & (0.5239)\\
 & stationary &  & $-1.3590$ & $-1.7775$ & $-2.3874$ & $-3.2747$\\
 & &  & (0.6500) & (0.5883) & (0.5490) & (0.5261)\\[0.4ex]
$\hat\tau_\mu$ & zero & $-1.5330$ & $-1.8373$ & $-2.0960$ & $-2.5924$ & $-3.4084$\\
 & & (0.8403) & (0.6949) & (0.6381) & (0.5836) & (0.5458)\\
 & stationary &  & $-1.8347$ & $-2.1226$ & $-2.6279$ & $-3.4404$\\
 & &  & (0.7100) & (0.6424) & (0.5853) & (0.5473)\\[0.4ex]
$\hat\tau_\tau$ & zero & $-2.1814$ & $-2.2927$ & $-2.4738$ & $-2.8644$ & $-3.5905$\\
 & & (0.7499) & (0.7086) & (0.6640) & (0.6110) & (0.5652)\\
 & stationary &  & $-2.2862$ & $-2.4731$ & $-2.8779$ & $-3.6118$\\
 & &  & (0.7133) & (0.6684) & (0.6129) & (0.5664)\\
\bottomrule
\end{tabular}
\end{center}
{\small Mean and, in parentheses, standard deviation of the limit of the Dickey--Fuller $t$-ratio ($c_0=0$) under $\phi=1+c/n$, no drift, from Proposition~\ref{prop:mom}; the stationary start requires $c<0$.  At $c=0$ the values are those of \cite{Nabeya1999}.}
\end{table}

\section{Application: the two Dickey--Fuller tests under local alternatives}
\label{sec:app}

The coefficient statistic $n(\hat\phi-\phi_0)$ of \cite{DickeyFuller1979} has, in the same three regressions and under the same local alternatives, the limit $(U-c_0V)/V$, whose distribution function at $x$ is $\PP\{U-(c_0+x)V\le0\}$: the distribution function of Theorem~\ref{th:single} at zero, with $c_0+x$ in place of $c_0$, which is the probability that a quadratic form in the path is negative, a one-dimensional inversion (the Supplementary Material).  The local power of the two tests can therefore be set side by side from one set of programs, exactly.  The comparison has been made by simulation \citep{Phillips1987,Stock1994,ElliottRothenbergStock1996}, and \cite{NabeyaTanaka1990b} computed the limiting power of the coefficient test, and of other ratios of quadratic forms, by Fredholm determinants and numerical inversion; the $t$-ratio, a ratio with a square root, was out of their reach.

Figure~\ref{fig:power} shows the power of the $5\%$ tests against $c$, with the critical values at the unit root with a zero start ($-1.94087$, $-2.86157$ and $-3.41045$ for the $t$-ratios, $-8.039$, $-14.094$ and $-21.711$ for the coefficient statistics, which the tables of \cite{Fuller1996} give to one decimal), with a zero and with a stationary start, together with the power envelope of \cite{ElliottRothenbergStock1996} in cases \Fk\ and \Fu, where, as they note, it is one curve: the point-optimal test against $\bar c$ rejects for large values of the log likelihood ratio $\bar cU-\frac12\bar c^2V$ of the zero-start process, a linear combination of $U$ and $V$, whose distribution function at $c=0$ and at $c=\bar c$ follows from Proposition~\ref{prop:mgf} by the inversion formula of \cite{GilPelaez1951}; its power at $\bar c=-7$, where they place the point of power one half, is $0.499$.

The GLS-detrended tests of \cite{ElliottRothenbergStock1996} are within reach as well.  With a constant their limit is that of case \Fk, since the quasi-demeaning is asymptotically negligible, so that the first panel of Figure~\ref{fig:power} is also their power.  With a trend the detrended path is $J(r)-r\{\lambda J(1)+3(1-\lambda)K\}$, $\lambda=(1-\bar c)/(1-\bar c+\bar c^2/3)$ at $\bar c=-13.5$, so that the numerator and the denominator of the statistic are quadratic forms in the Gaussian functionals of case \Ft, the numerator of rank one as in case \Fk, and Lemma~\ref{lem:roots}, the Girsanov forms and Theorem~\ref{th:single} apply with $D$ built from the same $\bSig$; the programs carry this case (\texttt{DG}).  The $5\%$ critical value of the $t$-ratio is $-2.8456$, so that the $-2.89$ of their Table~I has size $0.0446$; the coefficient form's is $-16.59$.  The fourth panel shows the two tests with the trend envelope, from the same inversion with the detrending at the alternative itself.

Four things stand out.  Without deterministic terms the two tests are indistinguishable and essentially efficient: both reach power one half within $0.1$ of $c=-7.0$, where the envelope does, and neither departs from it by more than $0.006$ at any $c$.  With a constant or a trend the coefficient test is the more powerful by a wide margin, at power one half by $c=-10.6$ against $-13.5$ in case \Fu\ and by $-15.9$ against $-17.8$ in case \Ft, and at $c=-10$ its power is $0.46$ against $0.30$ and $0.23$ against $0.18$; both are far below the envelope, whose point-optimal tests use the GLS-detrended data.  The stationary start moves the two tests in opposite directions: it raises the power of the $t$-ratio (one half at $c=-6.5$, $-13.0$ and $-17.5$) and lowers that of the coefficient test ($-7.1$, $-10.8$ and $-16.0$), so that in case \Fk\ the $t$-ratio with a stationary start lies above the envelope of the zero-start problem, whose optimal test is not optimal under a stationary start \citep{Elliott1999}.  And the GLS-detrended tests of the fourth panel lie on the trend envelope, the $t$-ratio and the coefficient form alike, with power one half at $c=-13.3$ against $-17.8$ and $-15.9$ for the least squares detrended tests, which is the gain that \cite{ElliottRothenbergStock1996} found by simulation; at their $\bar c=-13.5$ the envelope is $0.514$.

\begin{figure}[tbp]
\pgfplotsset{
  powerpanel/.style={width=0.48\textwidth, height=0.30\textwidth, xmin=-30, xmax=0, ymin=0, ymax=1, x dir=reverse,
    xtick={0,-10,-20,-30}, ytick={0,0.25,0.5,0.75,1}, yticklabels={$0$,$0.25$,$0.5$,$0.75$,$1$},
    tick label style={font=\footnotesize}, label style={font=\footnotesize}, title style={font=\footnotesize, yshift=-0.8ex},
    axis line style={line width=0.4pt}, every tick/.style={line width=0.4pt}, tick align=outside, tick pos=left,
    every axis plot/.append style={line width=0.7pt, mark=none, smooth}, xlabel={$c$}, xlabel style={yshift=0.5ex},
    legend style={font=\footnotesize, draw=none, fill=none, column sep=0.6em, /tikz/every even column/.append style={column sep=0.6em}}, legend cell align=left}
}
\begin{center}
\begin{tikzpicture}
\begin{groupplot}[group style={group size=2 by 2, horizontal sep=1.0cm, vertical sep=1.3cm, y descriptions at=edge left, x descriptions at=edge bottom}, powerpanel, ylabel={power}]
\nextgroupplot[title={$\hat\tau$, $\hat\rho$ (and DF-GLS, constant)}, legend to name=powerlegend, legend columns=3]
\addplot[black!35, line width=1.8pt] table[x=c,y=env] {figures/power_K.dat};
\addplot[blue!70!black] table[x=c,y=tau0] {figures/power_K.dat};
\addplot[red!80!black] table[x=c,y=rho0] {figures/power_K.dat};
\addplot[blue!70!black, dashed] table[x=c,y=tauS] {figures/power_K.dat};
\addplot[red!80!black, dashed] table[x=c,y=rhoS] {figures/power_K.dat};
\legend{power envelope, $t$-ratio, coefficient statistic, $t$-ratio{,} stationary start, coefficient statistic{,} stationary start}
\nextgroupplot[title={$\hat\tau_\mu$, $\hat\rho_\mu$}]
\addplot[black!35, line width=1.8pt] table[x=c,y=env] {figures/power_U.dat};
\addplot[blue!70!black] table[x=c,y=tau0] {figures/power_U.dat};
\addplot[red!80!black] table[x=c,y=rho0] {figures/power_U.dat};
\addplot[blue!70!black, dashed] table[x=c,y=tauS] {figures/power_U.dat};
\addplot[red!80!black, dashed] table[x=c,y=rhoS] {figures/power_U.dat};
\nextgroupplot[title={$\hat\tau_\tau$, $\hat\rho_\tau$}]
\addplot[blue!70!black] table[x=c,y=tau0] {figures/power_T.dat};
\addplot[red!80!black] table[x=c,y=rho0] {figures/power_T.dat};
\addplot[blue!70!black, dashed] table[x=c,y=tauS] {figures/power_T.dat};
\addplot[red!80!black, dashed] table[x=c,y=rhoS] {figures/power_T.dat};
\nextgroupplot[title={DF-GLS, trend: $\hat\tau^{\rm GLS}_\tau$, $\hat\rho^{\rm GLS}_\tau$}]
\addplot[black!35, line width=1.8pt] table[x=c,y=env] {figures/power_G.dat};
\addplot[blue!70!black] table[x=c,y=tau0] {figures/power_G.dat};
\addplot[red!80!black] table[x=c,y=rho0] {figures/power_G.dat};
\end{groupplot}
\end{tikzpicture}\\[0.4ex]
\pgfplotslegendfromname{powerlegend}
\end{center}
\caption{Local power of the $5\%$ Dickey--Fuller tests}
\label{fig:power}
{\small $\PP(\text{reject})$ under $\phi=1+c/n$, no drift, for the $t$-ratio and the coefficient statistic centered at the unit root, with the limiting $5\%$ critical values of the zero start; solid: zero start, dashed: stationary start; gray: the power envelopes of \cite{ElliottRothenbergStock1996}, one curve for cases \Fk\ and \Fu\ and the trend envelope in the fourth panel.  The first panel is also the power of the DF-GLS tests with a constant, whose limit is that of case \Fk; the fourth panel is the DF-GLS tests with a trend, detrended at $\bar c=-13.5$, in $t$-ratio and coefficient form, with a zero start.  From Theorem~\ref{th:single} and Proposition~\ref{prop:mgf} (\texttt{julia/application/power.jl}).}
\end{figure}

\section{Application: confidence sets for the persistence parameter}
\label{sec:ci}

The $t$-ratio centered at $c_0$ is the Wald statistic for the hypothesis $c=c_0$, and $n(\hat\phi-1)-c_0$ is its coefficient form; under $c=c_0$ their limits are $(U-c_0V)/\sqrt V$ and $(U-c_0V)/V$ with $c=c_0$.  The set of $c_0$ at which the observed statistic lies between the $\alpha/2$ and $1-\alpha/2$ quantiles of its limit is therefore a confidence set for $c$ of asymptotic level $1-\alpha$, and the $c_0$ at which it sits at the median is a median-unbiased estimator \citep{Stock1991,Andrews1993}.  For the coefficient form the set is Stock's, the pivot being a shift of the unit root statistic; for the $t$-ratio it is not, since the observed pivot $\hat\tau(1-c_0/\hat\rho)$ is a line in $c_0$ through $(\hat\rho,0)$.  Stock's intervals, and those of \cite{ElliottStock2001} for the GLS statistics and of \cite{Hansen1999}, come from simulated tables or a bootstrap over a grid of $c$; here the belts are exact at every $c_0$, including the explosive side, and they come with the initial condition.

Figure~\ref{fig:ci} shows the $95\%$ belts of the two pivots in case \Fu\ with a zero and a stationary start, with the median curves that give the estimator, and, for case \Fu\ and for the trend case with least squares and GLS detrending, the probability that a false value $c_0$ is excluded from the $95\%$ set when the true value is $-5$: the power of the two-sided test of $c_0$, which measures the length of the sets.  The two forms differ at the two ends.  The coefficient form excludes values on the unit root side far more often: at $c=-5$ it excludes $c_0=0$ with probability $0.107$ against $0.061$ for the $t$-ratio, and $c_0=2$ with $0.27$ against $0.13$; at $c=-10$ it excludes the unit root with $0.30$ against $0.18$.  On the stationary side the $t$-ratio is marginally the better one ($0.226$ against $0.222$ at $c_0=-10$).  The stationary start shortens the sets on that side ($0.236$ at $c_0=-10$, $0.765$ at $-20$ against $0.741$) and lengthens them slightly toward the unit root.  With a trend the least squares sets are the longest ($0.160$ at $c_0=-10$ when $c=-5$), and GLS detrending restores the lengths of the constant case on the stationary side ($0.230$, $0.516$ at $-15$, $0.741$ at $-20$, within $0.01$ of case \Fu) but not toward the unit root ($0.060$ at $c_0=0$), the two GLS forms being indistinguishable.  For the example, $\hat\rho_\mu=-8$ and $\hat\tau_\mu=-2$, the $95\%$ sets are $[-13.6,3.3]$ by the coefficient form and $[-13.4,3.7]$ by the $t$-ratio, with the median-unbiased estimates $-4.0$ and $-4.2$; a stationary start leaves the sets within $0.05$ and moves the estimates to $-4.4$ and $-4.3$.  Each set takes a few seconds.

\begin{figure}[tbp]
\pgfplotsset{
  cipanel/.style={width=0.46\textwidth, height=0.30\textwidth, xmin=-30, xmax=4, x dir=reverse,
    xtick={4,0,-10,-20,-30}, tick label style={font=\footnotesize}, label style={font=\footnotesize}, title style={font=\footnotesize, yshift=-0.8ex},
    axis line style={line width=0.4pt}, every tick/.style={line width=0.4pt}, tick align=outside, tick pos=left,
    every axis plot/.append style={line width=0.7pt, mark=none, smooth}, xlabel={$c_0$}, xlabel style={yshift=0.5ex},
    legend style={font=\footnotesize, draw=none, fill=none, column sep=0.6em}, legend cell align=left}
}
\begin{center}
\begin{tikzpicture}
\begin{groupplot}[group style={group size=2 by 2, horizontal sep=1.9cm, vertical sep=1.2cm, x descriptions at=edge bottom}, cipanel]
\nextgroupplot[title={95\% belt of the coefficient statistic $\hat\rho_\mu$}, ylabel={$\hat\rho_\mu$}, ymin=-55, ymax=10]
\addplot[red!80!black] table[x=c0,y=rlo0] {figures/ci_belts_U.dat};
\addplot[red!80!black, dashed] table[x=c0,y=rmed0] {figures/ci_belts_U.dat};
\addplot[red!80!black] table[x=c0,y=rhi0] {figures/ci_belts_U.dat};
\addplot[red!80!black, dotted] table[x=c0,y=rloS] {figures/ci_belts_U.dat};
\addplot[red!80!black, dotted] table[x=c0,y=rhiS] {figures/ci_belts_U.dat};
\addplot[black!60, line width=0.5pt] coordinates {(4,-8) (-30,-8)};
\nextgroupplot[title={95\% belt of the $t$-ratio centered at $c_0$}, ylabel={$t$-ratio pivot}, ymin=-4, ymax=4]
\addplot[blue!70!black] table[x=c0,y=tlo0] {figures/ci_belts_U.dat};
\addplot[blue!70!black, dashed] table[x=c0,y=tmed0] {figures/ci_belts_U.dat};
\addplot[blue!70!black] table[x=c0,y=thi0] {figures/ci_belts_U.dat};
\addplot[blue!70!black, dotted] table[x=c0,y=tloS] {figures/ci_belts_U.dat};
\addplot[blue!70!black, dotted] table[x=c0,y=thiS] {figures/ci_belts_U.dat};
\addplot[black!60, line width=0.5pt] coordinates {(4,-3) (-30,5.5)};
\nextgroupplot[title={probability of excluding $c_0$ at $c=-5$, case \Fu}, ylabel={probability}, ymin=0, ymax=1, ytick={0,0.25,0.5,0.75,1}, yticklabels={$0$,$0.25$,$0.5$,$0.75$,$1$}, legend to name=cilegend, legend columns=2]
\addplot[red!80!black] table[x=c0,y=r0] {figures/ci_excl_c5.dat};
\addplot[blue!70!black] table[x=c0,y=t0] {figures/ci_excl_c5.dat};
\addplot[red!80!black, dotted] table[x=c0,y=rS] {figures/ci_excl_c5.dat};
\addplot[blue!70!black, dotted] table[x=c0,y=tS] {figures/ci_excl_c5.dat};
\legend{coefficient statistic, $t$-ratio}
\nextgroupplot[title={probability of excluding $c_0$ at $c=-5$, trend}, ylabel={probability}, ymin=0, ymax=1, ytick={0,0.25,0.5,0.75,1}, yticklabels={$0$,$0.25$,$0.5$,$0.75$,$1$}]
\addplot[red!80!black] table[x=c0,y=rT] {figures/ci_excl_c5.dat};
\addplot[blue!70!black] table[x=c0,y=tT] {figures/ci_excl_c5.dat};
\addplot[red!80!black, dotted] table[x=c0,y=rG] {figures/ci_excl_c5.dat};
\addplot[blue!70!black, dotted] table[x=c0,y=tG] {figures/ci_excl_c5.dat};
\end{groupplot}
\end{tikzpicture}\\[0.4ex]
\pgfplotslegendfromname{cilegend}
\end{center}
\caption{Confidence sets for $c$ by inverting the Wald pivots, case \Fu}
\label{fig:ci}
{\small Top: the $2.5\%$ and $97.5\%$ quantiles (solid; dotted: stationary start) and the median (dashed) of the limit of the pivot under $c=c_0$, against $c_0$, for the coefficient statistic, $c_0+(U-c_0V)/V$, and the $t$-ratio centered at $c_0$, whose observed pivot is $\hat\tau_\mu(1-c_0/\hat\rho_\mu)$; the gray line is the worked example, $\hat\rho_\mu=-8$ and $\hat\tau_\mu=-2$, whose confidence set is the range of $c_0$ over which it lies inside the belt and whose median-unbiased estimate is where it crosses the median.  Bottom: the probability that the $95\%$ set excludes $c_0$ when the true value is $c=-5$: case \Fu\ with a zero (solid) and a stationary (dotted) start, and the trend case with least squares (solid) and GLS (dotted) detrending.  From Theorem~\ref{th:single} (\texttt{julia/application/ci.jl}).}
\end{figure}

\section{Conclusion}
\label{sec:concl}

The limiting distribution of the Dickey--Fuller $t$-ratio is an integral of special functions along a line in the complex plane: the Fredholm determinant that generates the numerator and the denominator of the limit is a polynomial of degree at most two in the argument of the numerator, which makes one of the two inversions explicit, and what remains is one integral, evaluated to twelve digits in a fraction of a second.  The representation covers the three regressions, local alternatives on both sides of the unit root, the $t$-ratio centered at the unit root or at any local null value, and a zero, a fixed or a stationary initial observation; a fixed initial observation enters the moment generating function through a closed-form quadratic term, and the integrand through the density of a noncentral quadratic form.  The moments are one-dimensional integrals.  The critical values, local powers and moments that have been obtained by simulation since \cite{DickeyFuller1979} are now exact, and so are the power envelopes, the local power of the GLS-detrended tests and the confidence sets for the persistence parameter.

The structure exploited here, a quadratic form of low rank in the numerator under a tilt of the denominator, arises wherever a $t$-ratio is formed from a Gaussian process: the limits of the Dickey--Fuller tests with a drift or a trend in the data-generating process \citep{NabeyaSorensen1994}, which Proposition~\ref{prop:fixed} already covers through the parameter $f$, and, after a convolution with a normal, the $t$-ratio in a predictive regression with a near-integrated regressor; Theorem~\ref{th:inv} applies to them without change.

\section*{Declaration of the Use of Generative AI}
An AI assistant (Claude, Anthropic) was used in drafting and editing the text and in writing and checking the replication code; all results were verified by the authors, who take full responsibility for them.
\begin{appendix}
\section{Proofs}
\label{app:proofs}

\begin{proof}[Proof of Theorem~\ref{th:inv}]
Assume first that
\begin{equation}
\int_{-\infty}^{\infty}\int_{-\infty}^{\infty}|M(\sigma+\iu u,\tau+\iu w)|\dr{u}\dr{w}<\infty.
\label{eq:L1}
\end{equation}
Then, by the inversion theorem for the bivariate two-sided Laplace transform, $(X,Y)$ has a continuous density
\[
f(x,y)=\frac{1}{(2\pi\iu)^2}\int_{\sigma-\iu\infty}^{\sigma+\iu\infty}\int_{\tau-\iu\infty}^{\tau+\iu\infty}M(s,t)e^{-sx-ty}\dr{t}\dr{s},
\]
and $|f(x,y)|\le Ce^{-\sigma x-\tau y}$ with $C$ equal to $(2\pi)^{-2}$ times (\ref{eq:L1}).  The map $(p,v)\mapsto(pv+c_0v^2,v^2)$ takes $\mathbb{R}\times(0,\infty)$ one-to-one onto itself with Jacobian $2v^2$, and $pv+c_0v^2-c_0v^2\le q\sqrt{v^2}$ if and only if $p\le q$, so
\[
\PP(Q_0\le q)=\int_{-\infty}^q\int_0^\infty2v^2f(pv+c_0v^2,v^2)\dr{v}\dr{p}.
\]
The inner integrand is bounded by $2Cv^2e^{-\sigma pv-(\tau+c_0\sigma)v^2}$, which is integrable locally uniformly in $p$ since $\tau+c_0\sigma>0$; hence the inner integral is continuous in $p$ and is the density of $Q_0$.  Inserting the inversion formula, the absolute integrand is $2v^2|M(s,t)|e^{-\sigma qv-(\tau+c_0\sigma)v^2}$, whose integral is $2\Xi_2(\sigma q,\tau+c_0\sigma)$ times (\ref{eq:L1}), so by Fubini the $v$-integral may be done first; it is $2\Xi_2(qs,t+c_0s)$, which gives (\ref{eq:dinv}).  For (\ref{eq:cinv}), integrate (\ref{eq:dinv}) over $p>q$: since $\int_q^\infty\int_0^\infty v^2e^{-\sigma pv-(\tau+c_0\sigma)v^2}\dr{v}\dr{p}<\infty$ for $\sigma>0$, the order of integration may be changed, and
\[
\int_q^\infty2\Xi_2(ps,t+c_0s)\dr{p}=2\int_0^\infty v^2e^{-(t+c_0s)v^2}\int_q^\infty e^{-psv}\dr{p}\dr{v}=\frac{2\Xi_1(qs,t+c_0s)}{s}.
\]
For $\sigma<0$ integrate over $p<q$ instead, where $\int_{-\infty}^qe^{-psv}\dr{p}=-e^{-qsv}/(sv)$.

In general, let $Z$ be standard normal and independent of $(X,Y)$, and for $\epsilon>0$ let $Q_\epsilon=(X+\epsilon Z-c_0Y)/\sqrt Y$.  The moment generating function of $(X+\epsilon Z,Y)$ is $M(s,t)e^{\epsilon^2s^2/2}$, and $|e^{\epsilon^2s^2/2}|=e^{\epsilon^2(\sigma^2-u^2)/2}$, so it satisfies (\ref{eq:L1}) by (\ref{eq:A}), and by the first part $Q_\epsilon$ has the continuous density $f_\epsilon(q)$ given by (\ref{eq:dinv}) with $e^{\epsilon^2s^2/2}$ inserted.  Let $g(q)$ denote the right side of (\ref{eq:dinv}).  For $\epsilon\le1$ the integrands of $f_\epsilon$ and $g$ are dominated by $e^{\sigma^2/2}$ times the integrand of (\ref{eq:B}), so $g$ is continuous on $I$ and, by dominated convergence, $f_\epsilon\to g$ uniformly on compact subsets of $I$ as $\epsilon\to0$.  The passage from $f_\epsilon$ to the density of $Q_0$ uses only this locally uniform convergence to a continuous limit, and is stated so that the proof of Theorem~\ref{th:single} can use it: \emph{if the densities $f_\epsilon$ of $Q_\epsilon$ converge as $\epsilon\to0$, locally uniformly on $I$, to a continuous function $g$, then the distribution function of $Q_0$ is continuous on $I$ and $g$ is its density there.}  Since $Q_\epsilon\to Q_0$ almost surely, hence in distribution, $\PP(Q_\epsilon\le q)\to\PP(Q_0\le q)$ at every continuity point $q$ of the distribution function of $Q_0$ \citep[Theorem~25.2]{Billingsley1995}.  So for $a<b$ in $I$ that are continuity points,
\[
\PP(a<Q_0\le b)=\lim_{\epsilon\to0}\PP(a<Q_\epsilon\le b)=\lim_{\epsilon\to0}\int_a^bf_\epsilon(q)\dr{q}=\int_a^bg(q)\dr{q},
\]
the last step by the uniform convergence on $[a,b]$.  The discontinuity points are countable, so every $a\in I$ lies between continuity points $a'<a<a''$ as close to $a$ as desired, and $\PP(Q_0=a)\le\int_{a'}^{a''}g(q)\dr{q}\to0$, $g$ being bounded on compact subsets of $I$: the distribution function of $Q_0$ is continuous on $I$, the identity extends to all $a<b$ in $I$, and $g$ is the density of $Q_0$ on $I$.  For (\ref{eq:cinv}), the first part gives $\PP(Q_\epsilon>q)$ as the integral with $e^{\epsilon^2s^2/2}$ inserted, which converges to the right side of (\ref{eq:cinv}) by dominated convergence, while $\PP(Q_\epsilon>q)\to\PP(Q_0>q)$ by the same theorem, $q$ being a continuity point; the case $\sigma<0$ is the same.
\end{proof}\begin{proof}[Proof of Lemma~\ref{lem:roots}]
Write the Gaussian path as $\xi$ with covariance operator $\Sigma$, $U=\frac12\langle\xi,A\xi\rangle-\frac12$ and $V=\langle\xi,B\xi\rangle$ with $B\ge0$ of trace class, and $A=R^*A_1R$ with $R$ of rank $r$ onto $\mathbb{R}^r$ and $A_1$ a nonsingular symmetric $r\times r$ matrix: $r=1$ with $A_1=1$ in case \Fk, where $U=\{J(1)^2-1\}/2$, and $r=2$ with $A_1$ indefinite in cases \Fu\ and \Ft, where $U+\frac12$ is $\frac12J(1)^2-\bar JJ(1)$ and $\frac12J(1)^2+2J(1)m_0-6J(1)m_1-6m_0^2+12m_0m_1$, $m_k=\int_0^1r^kJ\dr{r}$, whose matrices have determinants $-1$ and $0$ with rank $2$.  By the formula for the moment generating function of Gaussian quadratic forms, on the tube $M(s,t)=e^{-s/2}\det(I-2t\Sigma B-2s\Sigma A)^{-1/2}$, and since $(I-2t\Sigma B)^{-1}\Sigma=(\Sigma^{-1}-2tB)^{-1}=:\Sigma_t$ and $\det(I-2s\Sigma_tR^*A_1R)=\det(I_r-2sA_1R\Sigma_tR^*)$,
\[
D(s,t)=\det(I-2t\Sigma B)\,\det(I_r-2sA_1R\Sigma_tR^*),
\]
a polynomial in $s$ of degree $r$ with leading coefficient $(-2)^r\det A_1\det(R\Sigma_tR^*)\det(I-2t\Sigma B)$.  For $\Rea t<t_*$ the operator $\Sigma^{-1}-2tB$ has positive definite real part, hence so has its inverse $\Sigma_t$ and the $r\times r$ matrix $R\Sigma_tR^*$, which is therefore nonsingular (a complex symmetric matrix $G$ with $v^*(\Rea G)v>0$ for $v\ne0$ has $v^*Gv\ne0$); this is the accretivity argument of \cite{BaoKan2026}.  And $\det(I-2t\Sigma B)=D(0,t)\ne0$ for $\Rea t<t_*$, since $|M_V(t)|\le M_V(\Rea t)<\infty$.  So the degree is exactly $r$ and the roots are continuous in $t$.  At real $t=\tau$, $\Sigma_\tau$ is a covariance, $A_1R\Sigma_\tau R^*$ is similar to the symmetric matrix $(R\Sigma_\tau R^*)^{1/2}A_1(R\Sigma_\tau R^*)^{1/2}$, whose inertia is that of $A_1$ by Sylvester's law, and the roots $s_j=1/(2\lambda_j)$, $\lambda_j$ its eigenvalues, are real with the signs stated; $D(\cdot,\tau)$ is positive at $s=0$, vanishes only at the roots, and the component of $\{D>0\}$ containing $0$ is the interval of convergence of $M(\cdot,\tau)$.  Finally, $|M(s,t)|\le M(\sigma,\tau)<\infty$ on the tube, so $D(s,t)\ne0$ there: no root crosses the line $\Rea s=\sigma$ as $t$ moves along the contour, and the configuration at $t=\tau$ persists.  In the stationary cases the path is $(\gamma_0,W)$ with $\gamma_0$ an independent Gaussian coordinate, and (\ref{eq:U}) is again a quadratic form of rank two, in $(J(0),J(1))$ in case \Sk; the argument is unchanged.  With a fixed start the same argument, with the mean of $\xi$ carried along, shows that the weights do not depend on the start and that the configuration of the roots is the same.
\end{proof}\begin{proof}[Proof of Theorem~\ref{th:single} and Corollary~\ref{cor:closed}]
Theorem~\ref{th:inv} is applied to $(X,Y)=(U,V)$, on the tube $\Rea s=\sigma$, $\Rea t=\tau$ of Lemma~\ref{lem:roots}, through the smoothing in its proof rather than through hypothesis (\ref{eq:B}), which fails in two of the configurations the theorem covers: at $q=0$ when $c_0=0$, where $\Xi_2(0,t)$ does not depend on $s$ and $|M(s,t)|$ decays only like $|s|^{-r/2}$; and in case \Fk\ when $c_0\ne0$ and the parabola $X(v)$ has its vertex at a positive $v$, that is $c_0x<0$, where $2\Xi_2(xs,t+c_0s)$, by stationary phase at $v=-x/(2c_0)$, decays only like $|s|^{-1/2}$, as does $|M|$.  Hypothesis (\ref{eq:A}) holds: $|M(s,t)|\le C|P(t)|(1+|s|)^{-r/2}$ on the tube, because the roots $s_j(t)$, which are of the order of $|t|^{1/2}$ for large $|t|$, stay off the line $\Rea s=\sigma$, and $|P(t)|\le C\exp(-\frac12|\Ima t|^{1/2})$.  So, for $\epsilon>0$, $M(s,t)e^{\epsilon^2s^2/2}$ satisfies (\ref{eq:L1}), and by the first part of the proof of Theorem~\ref{th:inv} the variable $Q_\epsilon=(U+\epsilon Z-c_0V)/\sqrt V$ has the continuous density $f_\epsilon$ given by (\ref{eq:dinv}) with $e^{\epsilon^2s^2/2}$ inserted.  In that absolutely convergent integral the $s$- and $v$-integrals may be interchanged, by Fubini's theorem, once $\Xi_2$ is written out:
\begin{align*}
&\frac{1}{2\pi\iu}\int_{\sigma-\iu\infty}^{\sigma+\iu\infty}M(s,t)e^{\epsilon^2s^2/2}\,2\Xi_2(xs,t+c_0s)\dr{s}\\
&\qquad=\int_0^\infty2v^2e^{-tv^2}\Big\{\frac{1}{2\pi\iu}\int_{\sigma-\iu\infty}^{\sigma+\iu\infty}M(s,t)e^{\epsilon^2s^2/2}e^{-s(xv+c_0v^2)}\dr{s}\Big\}\dr{v},
\end{align*}
and the inner integral is $P(t)\varphi_{t,\epsilon}(X(v))$ with $\varphi_{t,\epsilon}=\varphi_t*g_\epsilon$, the convolution of (\ref{eq:varphi}) with the $N(0,\epsilon^2)$ density $g_\epsilon$, since $e^{\epsilon^2s^2/2}$ is the transform of $g_\epsilon$ and (\ref{eq:phit}) that of $\varphi_t$.  Hence
\[
f_\epsilon(x)=\frac{1}{2\pi\iu}\int_{\tau-\iu\infty}^{\tau+\iu\infty}P(t)\,\Phi^f_{x,\epsilon}(t)\dr{t},
\]
with $\Phi^f_{x,\epsilon}$ the first transform of (\ref{eq:Phidef}) with $\varphi_{t,\epsilon}$ in place of $\varphi_t$.  Now let $\epsilon\to0$.  Off $X=0$ the function $\varphi_t$ is continuous, so $\varphi_{t,\epsilon}(X)\to\varphi_t(X)$ for every $X\ne0$.  On the tube $|\varphi_t(X)|\le C(t)e^{-\sigma X}h(X)$, with $h(X)=X^{-1/2}\mathbf{1}\{X>0\}$ in case \Fk\ and $h(X)=1+|\log|X||$ otherwise, since $\Rea s_2<\sigma<\Rea s_1$ and $|K_0(z)|\le C(1+|\log|z||)e^{-\Rea z}$ for $\Rea z\ge0$, with $C(t)$ bounded by a power of $|t|$; and for $\epsilon\le1$ the smoothed function obeys the same bound with $h+1$ in place of $h$, because the convolution of $h$ with a positive radially decreasing kernel of unit mass is at most the Hardy--Littlewood maximal function of $h$, which for these two $h$ is a constant multiple of $h$, and the factor $e^{\sigma y}$ that the weight $e^{-\sigma X}$ leaves under the convolution is harmless when $\epsilon\le1$.  The dominating function $2v^2e^{-\tau v^2}e^{-\sigma X(v)}\{h(X(v))+1\}$ is integrable over $v>0$, as $\tau+c_0\sigma>0$ and the zeros of $X(v)$ are simple, except in case \Fk\ with $c_0>0$ at $x=-\sqrt{2c_0}$, where $X$ has a double root and $h(X(v))$ is not integrable, which is the point at which the density is infinite.  By dominated convergence $\Phi^f_{x,\epsilon}(t)\to\Phi^f_x(t)$, and then, as $|P(t)|C(t)$ is integrable along the line, $f_\epsilon(x)$ converges to the right side of (\ref{eq:single}), uniformly on compact sets of $x$ that avoid $-\sqrt{2c_0}$ in that case, and the limit is continuous there by the same domination.  The step in italics in the proof of Theorem~\ref{th:inv} then shows that the right side of (\ref{eq:single}) is the density of the limit and that its distribution function is continuous, on $(-\infty,-\sqrt{2c_0})$ and $(-\sqrt{2c_0},\infty)$ in that case, hence also at $-\sqrt{2c_0}$, where the singularity of the density is logarithmic and integrable, and on $\mathbb{R}$ otherwise.  For (\ref{eq:cinv}) the same steps with
\[
\frac{2\Xi_1(xs,t+c_0s)}{s}=2\int_0^\infty ve^{-tv^2}\,\frac{e^{-s(xv+c_0v^2)}}{s}\dr{v},\qquad
\frac{e^{-su}}{s}=\int_u^\infty e^{-sy}\dr{y},
\]
give $\PP(Q_\epsilon>x)$ as the integral over $v$ of $2tve^{-tv^2}$ times $P(t)\Psi_{t,\epsilon}(X(v))$, divided by $t$, with $\Psi_{t,\epsilon}=\Psi_t*g_\epsilon$, and as $\epsilon\to0$ this converges to the first form of $\Phi_x$ in (\ref{eq:Phidef}), $\Psi_t$ being continuous and bounded by $C(t)e^{-\sigma X}$ on the tube, while $\PP(Q_\epsilon>x)\to\PP(\tau>x)$ at every $x$ because the distribution function is continuous; the $1/t$ comes from writing $2ve^{-tv^2}$ as $2tve^{-tv^2}/t$.  The second form is the integration by parts
\[
\int_0^\infty\frac{\dr{}}{\dr{v}}\big(-e^{-tv^2}\big)\Psi_t\big(X(v)\big)\dr{v}=\Psi_t\big(\tfrac12\big)+\int_0^\infty e^{-tv^2}X'(v)\varphi_t\big(X(v)\big)\dr{v}.
\]  At $c_0=0$ the closed forms of Corollary~\ref{cor:closed} follow by writing the integrals over $(0,\frac12)$ as the difference of the integrals over $(0,\infty)$ and $(\frac12,\infty)$ in case \Fk, where $\int_0^\infty X^{k-1/2}e^{-\beta X^2-\gamma X}\dr{X}$ is $\cP_k$ and $\int_{1/2}^\infty$ is $\cQ_k$ after $X=\frac12(1+u)$, and by the definition (\ref{eq:G}) otherwise.
\end{proof}
\begin{lemma}[The tilted process given its start]
\label{lem:riccati}
In case \Fk\ with a fixed start $\gamma_0$, for real $t<t_*$ and real $u$,
\begin{equation}
\E\big\{e^{\iu uJ(1)+tV}\big\}=\{e^cF_0(z)\}^{-1/2}\exp\Big(-\frac{u^2}{2}\,\frac{F_1(z)}{F_0(z)}
+\iu u\,\frac{\gamma_0}{F_0(z)}+\gamma_0^2\,\frac{tF_1(z)}{F_0(z)}\Big),
\label{eq:riccati}
\end{equation}
with $F_0$, $F_1$ of case \Fk\ at $z=2t-c^2$.  Under the tilt $e^{tV}$, $J(1)$ given $J(0)=\gamma_0$ is therefore Gaussian with mean $\gamma_0/F_0$ and variance $F_1/F_0$, and the tilt multiplies the moment generating function of $V$ by $e^{\gamma_0^2tF_1/F_0}$.
\end{lemma}

\begin{proof}
Let $h(r,x)=\E\{\exp(\iu uJ(1)+t\int_r^1J(q)^2\dr{q})\mid J(r)=x\}$.  By the Feynman--Kac formula, $\partial_rh+cx\,\partial_xh+\frac12\partial_{xx}h+tx^2h=0$ with $h(1,x)=e^{\iu ux}$.  The trial $h=\exp\{\alpha(r)x^2+\beta(r)x+\delta(r)\}$ gives the Riccati system $\alpha'+2c\alpha+2\alpha^2+t=0$, $\beta'+(c+2\alpha)\beta=0$, $\delta'+\alpha+\frac12\beta^2=0$ with $\alpha(1)=0$, $\beta(1)=\iu u$, $\delta(1)=0$.  With $\alpha=k'/(2k)$ the first equation is $k''+2ck'+2tk=0$, solved by $k(r)=e^{-cr}g(r)$, $g(r)=\cosh\mu(1-r)-(c/\mu)\sinh\mu(1-r)$, $\mu^2=c^2-2t=-z$, which satisfies $k'(1)=0$; then $g(0)=C-cS=F_0$, $k'(0)/k(0)=2tS/F_0$ and $\alpha(0)=tF_1/F_0$.  The second gives $\beta=\iu ue^{-cr}/k(r)$, so $\beta(0)=\iu u/F_0$.  The third gives $\delta(0)=\int_0^1\{k'/(2k)+\frac12\beta^2\}\dr{r}=-\frac12\log\{e^cF_0\}-\frac12u^2\int_0^1g(r)^{-2}\dr{r}$, and since $g''=\mu^2g$, the Wronskian of $g$ and $g_2(r)=\sinh\mu(1-r)/\mu$ is constant, $gg_2'-g'g_2=-1$, so that $\int_0^1g^{-2}\dr{r}=g_2(0)/g(0)=S/F_0$.  Evaluating $h$ at $(0,\gamma_0)$ gives (\ref{eq:riccati}); its continuation to complex $t$ and $u$ holds on the tube.
\end{proof}\begin{proof}[Proof of Proposition~\ref{prop:fixed}]
In case \Fk, $R=\frac12\{J(1)^2-\gamma_0^2\}$, and by Lemma~\ref{lem:riccati}, under the tilt, $J(1)\sim N(\gamma_0/F_0,F_1/F_0)$ while $M_V$ gains $e^{\gamma_0^2tF_1/F_0}$; the moment generating function of $\frac12X^2$ for $X\sim N(m,w)$ is $(1-sw)^{-1/2}\exp\{\frac12sm^2/(1-sw)\}$, which gives (\ref{eq:QK}) after the subtraction of the zero-start value.  In cases \Fu\ and \Ft\ the dependence (\ref{eq:fixedstart})--(\ref{eq:g}) of $\log M$ on the start is Theorems~1 and 2 of \cite{NabeyaSorensen1994}, in the notation of Proposition~\ref{prop:mgf}, whose parameter $f$ is $c\gamma_0$ and $c^2\gamma_0$ here.  For the partial fractions, note that $R$ is invariant to $J\mapsto J+a+br$ (to $J\mapsto J+a$ in case \Fu), as the regression is; one checks this directly on (\ref{eq:U}).  The functionals $(J(0),J(1),\bar J,K)$ change by $(a,a+b,a+b/2,a/2+b/3)$, so $R$ is a quadratic form in two linear combinations of them that are invariant, whose coefficients do not involve the start.  With $J=\gamma_0h+J_0$, $h(r)=e^{cr}$ and $J_0$ the zero-start process, these combinations are Gaussian with means proportional to $\gamma_0$, also under the tilt, and the covariance under the tilt does not involve $\gamma_0$; so $R=\frac12(X_1^2-X_2^2)$ with $X_1$, $X_2$ independent Gaussian in the eigenbasis of the quadratic form, of variances $2\kappa_1$, $-2\kappa_2$ with the weights of the zero start, and means proportional to $\gamma_0$.  The moment generating function of this noncentral form is $\prod_{k=1}^2(1-2s\kappa_k)^{-1/2}\exp\{\frac s2\sum_{k=1}^2(-1)^{k-1}m_k^2/(1-2s\kappa_k)\}$, so that $Q(s,t)-Q(0,t)$ is of the form (\ref{eq:partial}) with $m_k^2=\gamma_0^2\mu_k$, and the $\mu_k$ are read off the residues of (\ref{eq:g}) at the roots: with $D(s,t)=e^cF_2(s-s_1)(s-s_2)$ the residue of $Q$ at $s_k$ is $c^{2k}(t+\frac12s_k^2+cs_k)g(z)/\{F_2(s_k-s_{k'})\}$, $k'\ne k$, and that of the right side of (\ref{eq:partial}) is $-\frac12\mu_1s_1^2$ at $s_1$ and $\frac12\mu_2s_2^2$ at $s_2$.  Both sides are rational in $s$ with the same poles and the same value at $s=0$, and tend to constants as $s\to\infty$, so they agree once the residues do.
\end{proof}\begin{proof}[Proof of Corollary~\ref{cor:fixed}]
The proof of Theorem~\ref{th:single} goes through with $M_{\gamma_0}=M_0e^{\gamma_0^2Q}$ in place of $M_0$ once two of its ingredients are checked: hypothesis (\ref{eq:A}), and the bounds on $\varphi_t$ and $\Psi_t$ that carry the passage $\epsilon\to0$.

For (\ref{eq:A}) we show that for every real $(\sigma,\tau)$ in the domain of $M_0$ and every $(s,t)$ on the tube $\Rea s=\sigma$, $\Rea t=\tau$,
\begin{equation}
\Rea Q(s,t)\le Q(\sigma,\tau).
\label{eq:ReQ}
\end{equation}
Under the tilt by $e^{\sigma U+\tau V}$ the process $J=\gamma_0h+J_0$ is Gaussian, with the covariance of the tilted zero-start process, which does not involve $\gamma_0$, and a mean proportional to $\gamma_0$ (the proof of Proposition~\ref{prop:fixed}), and $uU+wV$, with $u=\Ima s$ and $w=\Ima t$, is a quadratic form in $J$ plus a constant.  In the Karhunen--Lo\`eve coordinates of that covariance the form is $\sum_j\beta_j(\xi_j+m_j)^2$ with $\xi_j$ independent standard normal, real $\beta_j$ with $\sum_j|\beta_j|<\infty$, and $m_j$ proportional to $\gamma_0$, so that its characteristic function at the argument $1$ is $\prod_j(1-2\iu\beta_j)^{-1/2}\exp\{\iu\beta_jm_j^2/(1-2\iu\beta_j)\}$, and since $\Rea\{\iu\beta_j/(1-2\iu\beta_j)\}=-2\beta_j^2/(1+4\beta_j^2)\le0$, its modulus is at most that of the central form, which is its value at $\gamma_0=0$.  In terms of the moment generating functions this reads $|M_{\gamma_0}(s,t)|/M_{\gamma_0}(\sigma,\tau)\le|M_0(s,t)|/M_0(\sigma,\tau)$, which is (\ref{eq:ReQ}).  Hence $|M_{\gamma_0}(s,t)|\le e^{\gamma_0^2Q(\sigma,\tau)}|M_0(s,t)|$ on the tube, and (\ref{eq:A}) for $M_{\gamma_0}$ follows from (\ref{eq:A}) for $M_0$.

The passage $\epsilon\to0$ used, on the tube, $|\varphi_t(X)|\le C(t)e^{-\sigma X}h(X)$ and $|\Psi_t(X)|\le C(t)e^{-\sigma X}$, with $C(t)$ bounded by a power of $|t|$, and the integrability over $v>0$ of the dominating function $2v^2e^{-\tau v^2}e^{-\sigma X(v)}\{h(X(v))+1\}$.  For (\ref{eq:phinc}) the bounds hold with $h(\tilde X)=\tilde X^{-1/2}\mathbf{1}\{\tilde X>0\}$, $\tilde X=X+\frac12\gamma_0^2$: by $|\cosh z|\le e^{|z|}$ and $|b|(2\tilde X)^{1/2}\le\eta\tilde X+|b|^2/(2\eta)$,
\[
\big|\cosh\{ms_1(2\tilde X)^{1/2}\}\big|\,e^{-\Rea s_1\tilde X}\le e^{|ms_1|^2/(2\eta)}\,e^{-(\Rea s_1-\eta)\tilde X}\qquad(\eta>0),
\]
where $|ms_1|=\gamma_0/|F_1|$ is bounded on the line and $\Rea s_1-\sigma$ is bounded away from zero there, so that $\eta$ can be fixed below it; and the prefactors $e^{\gamma_0^2tF_1/F_0}=e^{\gamma_0^2Q(0,t)}$ and $e^{-s_1m^2/2}$, with $s_1m^2=\gamma_0^2/(F_0F_1)$, are bounded on the line, the first by (\ref{eq:ReQ}) at $\sigma=0$.  For (\ref{eq:phinc2}) the same inequality bounds $\cosh(b_1l)$ and $\cosh\{b_2(2Y+l^2)^{1/2}\}$ against $e^{-al^2/2}$ and $e^{-s_dX}$, since $\Rea a=\Rea(s_1-s_2)$, $\Rea s_1-\sigma$ and $\sigma-\Rea s_2$ are bounded away from zero on the line and $|b_k|=|m_ks_k|$ is bounded there, by (\ref{eq:partial}): $s_1-s_2$ and $t+\frac12s_k^2+cs_k$ are of the orders $|t|^{1/2}$ and $|t|$, and $g(z)/F_2(z)$ of the order $|t|^{-1/2}$; what remains is $\int_0^\infty e^{-\Rea a\,l^2/4}(2|Y|+l^2)^{-1/2}\dr{l}$, the logarithmic singularity of $h(X)=1+|\log|X||$ at $X=0$, which is the kink of the zero start.  The prefactors $e^{\gamma_0^2Q(0,t)}$ and $e^{(s_2m_2^2-s_1m_1^2)/2}$ are bounded on the line, the first by (\ref{eq:ReQ}) and the second because $s_km_k^2=\gamma_0^2\mu_ks_k$ tends to zero.  With these bounds the dominated convergence of the proof of Theorem~\ref{th:single} applies as before: in cases \Fu\ and \Ft\ with the same $h$ and the same simple zeros of $X(v)$, and in case \Fk\ with $h(\tilde X(v))$, $\tilde X(v)=c_0v^2+xv+\frac12(1+\gamma_0^2)$, whose zeros are simple except for $c_0>0$ at $x=-\{2c_0(1+\gamma_0^2)\}^{1/2}$, where $\tilde X$ has a double root, $h(\tilde X(v))$ is not integrable, and the density is infinite; it is the largest value of $\tau=-\{\frac12(1+\gamma_0^2)+c_0V\}/\sqrt V$ on the event $J(1)=0$, where $R=-\frac12\gamma_0^2$.  Near the double root $\varphi_t$ behaves like $\tilde X^{-1/2}$, as in the zero-start case, and the singularity of the density is logarithmic.
\end{proof}\begin{proof}[Proof of (\ref{eq:Dstat}) and (\ref{eq:DstatK})]
Since $\gamma_0^2\sim\sigma_0^2\chi^2_1$ with $\sigma_0^2=-1/(2c)$, $\E\exp(\gamma_0^2Q)=(1-2\sigma_0^2Q)^{-1/2}=(1+Q/c)^{-1/2}$, and $D(1+Q/c)$ is (\ref{eq:Dstat}).  In case \Sk, by Lemma~\ref{lem:riccati}, averaging over $\gamma_0\sim N(0,\sigma_0^2)$ with the weight $e^{\gamma_0^2tF_1/F_0}$ leaves $J(0)$ Gaussian with variance $\tilde\sigma^2=\sigma_0^2/(1-2\sigma_0^2tF_1/F_0)$ and multiplies the moment generating function of $V$ by $(1-2\sigma_0^2tF_1/F_0)^{-1/2}$, so that under the tilt the pair $(J(0),J(1))$ is Gaussian with $\Var J(0)=\tilde\sigma^2$, $\Cov\{J(0),J(1)\}=\tilde\sigma^2/F_0$ and $\Var J(1)=\tilde\sigma^2/F_0^2+F_1/F_0$.  With $R=\frac12\{J(1)^2-J(0)^2\}$ and $\Lambda={\rm diag}(-1,1)$, $\E e^{sR}=\det(I-s\Sigma\Lambda)^{-1/2}$, so $D_{\rm S}=e^cF_0(1-2\sigma_0^2tF_1/F_0)\{1-s(\Sigma_{22}-\Sigma_{11})-s^2\det\Sigma\}$.  Here $\det\Sigma=\tilde\sigma^2F_1/F_0$, and $\Sigma_{22}-\Sigma_{11}$ vanishes: by $C^2-zS^2=1$, $1-2tS^2=(C-cS)(C+cS)=F_0(C+cS)$, so that $\tilde\sigma^2/F_0^2+F_1/F_0-\tilde\sigma^2$, which equals $\{\sigma_0^2(C+cS)+F_1-\sigma_0^2F_0\}\tilde\sigma^2/(\sigma_0^2F_0)$ after the substitution of $\tilde\sigma^2$, is $(2c\sigma_0^2S+S)\tilde\sigma^2/(\sigma_0^2F_0)=0$ since $2c\sigma_0^2=-1$.  Multiplying out gives (\ref{eq:DstatK}).
\end{proof}
\end{appendix}

\bibliographystyle{econsoc}
\bibliography{LTR_BBK}

\end{document}
